\documentclass[12pt,a4paper]{report}
\usepackage[margin=2.5cm]{geometry}
\usepackage{fontspec}
\setmainfont{lmroman12-regular.otf}[BoldFont=lmroman12-bold.otf,ItalicFont=lmroman12-italic.otf,BoldItalicFont=lmroman10-bolditalic.otf]
\usepackage{polyglossia}
\setdefaultlanguage{indonesian}
\usepackage{amsmath,amssymb,array,booktabs,longtable,graphicx}
\usepackage{tikz,pgfplots}
\usetikzlibrary{arrows.meta,positioning,shapes.geometric}
\pgfplotsset{compat=1.18}
\makeatletter\def\input@path{{bibliography/}}\makeatother
\usepackage[backend=biber,style=filkom-authoryear,language=english,autolang=none]{biblatex}
\DeclareLanguageMapping{indonesian}{english}
\addbibresource{bibliography/references.bib}
\usepackage{xurl,hyperref}
\hypersetup{hidelinks}
\emergencystretch=3em
\title{Dokumentasi Penelitian untuk P1\\Klasifikasi Arah XAUUSD}
\author{Nathan}
\date{8 Oktober 2026 -- Draf untuk pembahasan P1}
\begin{document}
\renewcommand{\chaptername}{BAB}
\renewcommand{\contentsname}{DAFTAR ISI}
\maketitle
\chapter*{Status dokumen}
Dokumen ini adalah salinan kurasi Bab 1--4 untuk pembahasan P1, bukan naskah yang telah disahkan. Naskah resmi tetap pada repositori penelitian privat. Hasil K3 terbatas pada snapshot retrospektif. Perhitungan manual dan latihan tersedia terpisah pada direktori belajar. Persentase kemajuan formal ditentukan oleh pembimbing terhadap target berita acara P0.
\tableofcontents
\chapter{PENDAHULUAN}

\section{Latar Belakang}
Emas merupakan salah satu aset yang banyak diperdagangkan di pasar keuangan global dan sering digunakan sebagai instrumen lindung nilai (\emph{hedge}) ketika ketidakpastian ekonomi meningkat. Dalam pasar valuta asing, emas diperdagangkan melalui XAUUSD, yaitu harga satu troy ounce emas dalam dolar Amerika Serikat. Pergerakan harganya berkaitan dengan berbagai kondisi ekonomi dan pasar. \textcite{dinh2022economic} menunjukkan hubungan volatilitas dan korelasi logam mulia dengan pasar saham, neraca perdagangan, kepercayaan konsumen, jumlah uang beredar, dan inflasi. \textcite{wang2023asymmetric} juga menemukan hubungan asimetris antara imbal hasil logam mulia, harga minyak mentah, dan nilai tukar dolar Amerika Serikat dengan mempertimbangkan volatilitas pasar serta suku bunga. Pada masa guncangan global, hubungan lintas pasar tersebut dapat menguat, sebagaimana ditunjukkan oleh \textcite{yousaf2022information} melalui transmisi informasi antara emas dan minyak mentah selama wabah COVID-19. Banyaknya faktor yang terlibat membuat perkiraan pergerakan harga emas menjadi persoalan yang penting sekaligus menantang.

Untuk berbagai keperluan analisis, informasi mengenai arah pergerakan harga berikutnya dapat lebih berguna daripada perkiraan nilai harga secara tepat. Oleh karena itu, sejumlah penelitian merumuskan prediksi harga sebagai masalah klasifikasi arah. \textcite{sadorsky2021predicting} memprediksi arah harga emas dan perak dengan model logit dan beberapa pengklasifikasi berbasis pohon menggunakan indikator teknikal. Model berbasis pohon menghasilkan kinerja yang lebih baik daripada model linear pada sebagian besar konfigurasi yang diuji. \textcite{basher2022forecasting} kemudian memasukkan suku bunga, inflasi, dan volatilitas pasar dalam prediksi arah harian emas dan Bitcoin. Kedua penelitian tersebut menunjukkan bahwa jenis model dan kelompok informasi yang digunakan dapat memengaruhi kinerja klasifikasi arah.

Informasi untuk memprediksi arah harga dapat berasal dari riwayat harga maupun kondisi di luar pasar. Indikator teknikal mengubah riwayat harga menjadi representasi momentum atau tren yang lebih ringkas. Sebagai contoh, \textcite{mamoudan2023hybrid} memanfaatkan turunan MACD, \emph{moving average}, dan Ichimoku untuk mengklasifikasikan sinyal teknikal pada pasar emas sebagai \emph{FAKE} atau \emph{REAL}. Di sisi lain, variabel makroekonomi membawa informasi mengenai kondisi ekonomi. \textcite{sobti2021how} menunjukkan bahwa pengumuman makroekonomi memengaruhi proses pembentukan harga emas sepanjang hari perdagangan. Temuan ini mendukung pengujian informasi kalender ekonomi pada resolusi intrahari, dengan tetap memperhatikan waktu ketika informasi tersebut benar-benar tersedia.

Penambahan fitur belum tentu selalu meningkatkan kinerja karena dapat menambah dimensi, derau, dan risiko pemodelan berlebih (\emph{overfitting}). \textcite{taneva2025framework} telah memadukan data finansial, ekonomi, dan sentimen untuk meramalkan harga emas harian, sedangkan \textcite{velez2026predictive} membandingkan 32 kombinasi model dan fitur untuk peramalan harga emas bulanan. Kedua studi tersebut memperlihatkan pentingnya penggabungan dan perbandingan fitur. Namun, dalam batas studi yang ditelaah, perubahan kinerja akibat penambahan kelompok harga, teknikal, dan makroekonomi secara bertahap belum diukur untuk klasifikasi tiga arah XAUUSD H1 pada observasi uji yang sama.

Resolusi satu jam (H1) dipilih untuk mengamati dinamika intrahari dengan jumlah pengamatan yang lebih banyak daripada data harian. Penggunaan data H1 untuk XAUUSD memiliki preseden pada penelitian \textcite{arif2026multivariate}, yang meramalkan harga penutupan dengan LSTM dan validasi deret waktu bergulir. Pada konteks berbeda, \textcite{hisyam2025implementasi} menggunakan XGBoost dan \emph{walk-forward validation} untuk meramalkan harga emas Antam harian. Penelitian ini menerapkan Regresi Logistik Multinomial sebagai pembanding linear dan XGBoost sebagai pembanding nonlinear pada setiap konfigurasi fitur. Kedua model membantu membaca perubahan kinerja setelah penambahan fitur; perbedaan skor antarmodel ditafsirkan secara deskriptif, bukan sebagai uji linearitas data.

Pengukuran perubahan kinerja juga harus mengikuti urutan waktu. Pembagian data deret waktu secara acak dapat menyebabkan kebocoran informasi masa depan (\emph{look-ahead bias}) dan menghasilkan penilaian yang tidak mewakili penggunaan pada data baru. Penelitian ini karena itu menggunakan evaluasi \emph{walk-forward} dengan jeda antara data pelatihan dan pengujian. Macro-F1 digunakan sebagai metrik utama karena distribusi kelas BUY, HOLD, dan SELL dapat tidak seimbang, sedangkan balanced accuracy digunakan sebagai metrik pendamping.

Berdasarkan uraian tersebut, penelitian ini mengukur perubahan kinerja klasifikasi arah XAUUSD H1 setelah fitur teknikal dan makroekonomi ditambahkan secara bertahap. Konfigurasi harga (K1), harga dan teknikal (K2), serta harga, teknikal, dan makroekonomi (K3) dibandingkan menggunakan Regresi Logistik Multinomial dan XGBoost pada periode uji kronologis yang sama. Perubahan $K2-K1$ dan $K3-K2$ dibaca dalam masing-masing model, sedangkan perbedaan respons kedua model dijelaskan secara deskriptif. Kontribusi fitur ditafsirkan secara kondisional terhadap data, model, periode, dan prosedur penyetelan yang digunakan. Konfigurasi fitur makroekonomi (K3) dievaluasi menggunakan arsip kalender beku MetaTrader 5, sedangkan keterbatasan \emph{snapshot} retrospektif tanpa arsip \emph{vintage} independen membatasi interpretasi kinerja untuk penerapan waktu nyata.

\section{Rumusan Masalah}
Berdasarkan latar belakang tersebut, rumusan masalah penelitian ini adalah sebagai berikut.
\begin{enumerate}
    \item Bagaimanakah kinerja klasifikasi arah XAUUSD H1 pada konfigurasi fitur harga (K1), harga dan teknikal (K2), serta harga, teknikal, dan makroekonomi (K3) ketika masing-masing diterapkan pada Regresi Logistik Multinomial dan XGBoost?
    \item Bagaimanakah besar dan konsistensi perubahan kinerja setelah penambahan fitur teknikal dan fitur makroekonomi apabila dievaluasi pada beberapa periode uji berurutan dan pada kedua model yang digunakan?
\end{enumerate}

\section{Tujuan}
Penelitian ini bertujuan untuk:
\begin{enumerate}
    \item mengukur kinerja klasifikasi arah XAUUSD H1 secara \emph{out-of-sample} untuk K1, K2, dan K3 pada Regresi Logistik Multinomial serta XGBoost; dan
    \item mengukur serta menganalisis besar dan konsistensi perubahan kinerja antarkonfigurasi pada beberapa periode uji berurutan menggunakan Regresi Logistik Multinomial dan XGBoost.
\end{enumerate}

\section{Manfaat}
Hasil penelitian ini diharapkan memberikan manfaat sebagai berikut.
\begin{enumerate}
    \item Memberikan bukti terukur bagi bidang komputasi cerdas mengenai perubahan kinerja klasifikasi arah XAUUSD H1 setelah penambahan fitur teknikal dan makroekonomi dalam batas data, model, dan periode pengujian.
    \item Menyediakan prosedur perbandingan konfigurasi fitur bertahap dan evaluasi kronologis yang dapat diadaptasi oleh penelitian selanjutnya.
    \item Memberikan gambaran bagi peneliti dan praktisi mengenai perubahan kinerja yang diperoleh ketika kelompok fitur teknikal dan makroekonomi ditambahkan secara bertahap.
\end{enumerate}

\section{Batasan Masalah}
Agar pembahasan tetap terarah, penelitian ini dibatasi sebagai berikut.
\begin{enumerate}
    \item Objek penelitian adalah XAUUSD dengan timeframe satu jam (H1), sehingga simpulan dibatasi pada instrumen dan timeframe tersebut.
    \item Target pemodelan terdiri atas tiga label arah, yaitu BUY, HOLD, dan SELL, yang dibentuk dari perubahan harga penutupan satu bar berikutnya. Label tersebut tidak dimaksudkan sebagai rekomendasi transaksi.
    \item Model yang digunakan adalah Regresi Logistik Multinomial dan XGBoost.
    \item Fitur makroekonomi berasal dari data terstruktur kalender ekonomi MetaTrader 5, meliputi nilai indikator dan penanda peristiwa. Teks berita dan analisis sentimen tidak termasuk dalam penelitian.
    \item Kinerja diukur menggunakan macro-F1 dan balanced accuracy melalui evaluasi \emph{walk-forward}.
    \item Analisis bersifat prediktif, tidak menyimpulkan hubungan sebab-akibat ekonomi, menjanjikan keuntungan transaksi, atau menghasilkan sistem eksekusi order otomatis.
\end{enumerate}

\section{Sistematika Pembahasan}
Skripsi ini disusun dalam enam bab sebagai berikut.
\begin{enumerate}
    \item BAB 1 PENDAHULUAN menjelaskan latar belakang, rumusan masalah, tujuan, manfaat, batasan masalah, dan sistematika pembahasan.
    \item BAB 2 LANDASAN KEPUSTAKAAN membahas penelitian terdahulu serta teori objek, fitur, model, dan evaluasi penelitian.
    \item BAB 3 METODOLOGI menjelaskan protokol pengumpulan dan audit data, target, konfigurasi fitur, pelatihan, serta evaluasi kronologis.
    \item BAB 4 HASIL menyajikan hasil audit, sampel evaluasi, kinerja model, perubahan berpasangan, dan analisis sensitivitas.
    \item BAB 5 PEMBAHASAN menghubungkan hasil dengan rumusan masalah, konsistensi kontribusi kelompok fitur, penelitian terdahulu, dan ancaman validitas.
    \item BAB 6 PENUTUP memuat kesimpulan serta saran berdasarkan hasil dan keterbatasan penelitian.
\end{enumerate}

\chapter{LANDASAN KEPUSTAKAAN}

Bab ini membahas penelitian terdahulu yang berkaitan dengan prediksi harga emas serta teori yang digunakan dalam penelitian. Kajian pustaka menunjukkan posisi penelitian berdasarkan objek, target, fitur, model, dan evaluasi. Landasan teori menjelaskan konsep-konsep yang mendasari pembentukan tiga konfigurasi fitur dan pengukuran perubahan kinerjanya.

\section{Kajian Pustaka}

Kajian penelitian mutakhir menggunakan artikel jurnal final tahun 2021--2026, sedangkan rujukan yang lebih lama digunakan sebagai dasar teori. Studi yang dibahas mencakup klasifikasi arah, hubungan makroekonomi, penggabungan fitur, model prediksi, dan validasi temporal.

\subsection{Klasifikasi Arah dan Representasi Teknikal}

\textcite{sadorsky2021predicting} mengklasifikasikan arah harian ETF emas GLD dan ETF perak SLV untuk horizon 1--20 hari menggunakan 13 indikator teknikal. Penelitian tersebut membandingkan model logit, bagging, boosting, dan random forest. Model berbasis pohon memberikan hasil yang lebih baik daripada model linear pada sebagian besar konfigurasi pengujian. Studi ini menjadi dasar penggunaan model dengan karakter linear dan berbasis pohon, sedangkan penelitian yang diusulkan menggunakan harga spot XAUUSD H1 dan target tiga kelas.

\textcite{basher2022forecasting} memprediksi arah harian emas dan Bitcoin dengan fitur teknikal dan makroekonomi, termasuk suku bunga, inflasi, dan volatilitas pasar. Model logit, bagging, dan random forest dievaluasi menggunakan validasi silang deret waktu serta analisis kepentingan variabel. Penelitian tersebut menunjukkan bahwa fitur makroekonomi dapat disertakan dalam klasifikasi arah dan dievaluasi secara temporal. Penelitian ini melanjutkan gagasan tersebut pada klasifikasi BUY, HOLD, dan SELL dalam timeframe satu jam.

\textcite{mamoudan2023hybrid} mengklasifikasikan sinyal teknikal sebagai \emph{FAKE} atau \emph{REAL}. Arsitekturnya menggabungkan CNN--BiGRU, pemilihan fitur dengan Moth-Flame Optimization, dan penyetelan dengan Firefly Algorithm. Kandidat masukannya mencakup MACD, \emph{moving average}, dan Ichimoku. Studi tersebut memperlihatkan pemanfaatan transformasi teknikal, tetapi targetnya berupa validitas sinyal dan bukan arah harga berikutnya. Dalam penelitian ini, indikator teknikal digunakan untuk merepresentasikan momentum dan tren dari riwayat harga agar dapat dipelajari oleh model linear maupun nonlinear.

\subsection{Hubungan Makroekonomi dan Informasi Intrahari}

Hubungan antara pasar emas dan kondisi ekonomi telah dikaji melalui beberapa pendekatan ekonometrika. \textcite{dinh2022economic} menggunakan GARCH-MIDAS dan DCC-MIDAS untuk menganalisis volatilitas dan korelasi logam mulia dalam kaitannya dengan variabel makroekonomi. \textcite{wang2023asymmetric} menggunakan regresi kuantil untuk menilai hubungan asimetris harga minyak dan dolar Amerika Serikat dengan imbal hasil logam mulia, disertai federal funds rate dan VIX sebagai variabel kontrol. Hasil kedua studi menunjukkan bahwa hubungan tersebut dapat berbeda menurut kondisi pasar dan bagian distribusi yang diamati. Oleh karena itu, manfaat fitur makroekonomi untuk prediksi perlu diukur pada tugas dan periode pengujian yang ditetapkan.

Pada resolusi intrahari, \textcite{sobti2021how} menganalisis data satu menit periode 2013--2018 dan 62 pengumuman makroekonomi untuk mempelajari pembentukan harga di pasar emas. Temuan tersebut menunjukkan bahwa informasi pengumuman ekonomi relevan sepanjang hari perdagangan. \textcite{yousaf2022information} juga menggunakan data per jam untuk mengkaji limpahan informasi dan efektivitas lindung nilai pasangan minyak--emas serta minyak--Bitcoin selama COVID-19. Kedua studi ini mendukung penggunaan resolusi intrahari dan pengujian informasi eksternal, meskipun kontribusinya terhadap klasifikasi XAUUSD H1 tetap harus diukur secara \emph{out-of-sample}.

\subsection{Penggabungan Fitur, Model, dan Validasi Temporal}

\textcite{taneva2025framework} memadukan data finansial, makroekonomi, dan sentimen dengan metode ekonometrika, pembelajaran mesin, dan pembelajaran mendalam untuk meramalkan harga emas harian periode 2014--2024. Penelitian tersebut menggunakan target harga kontinu dan metrik regresi. Selain perbedaan target, penggunaan fitur berfrekuensi rendah pada keputusan H1 memerlukan pemeriksaan waktu ketersediaan data agar informasi masa depan tidak masuk ke fitur.

\textcite{velez2026predictive} meramalkan harga emas kontinu bulanan periode Januari 2020--Juni 2026 melalui 32 kombinasi model dan fitur dengan validasi temporal. Studi tersebut memberikan preseden langsung bagi perbandingan konfigurasi fitur. Penelitian ini menggunakan rancangan yang berbeda, yaitu tiga konfigurasi harga--teknikal--makro yang ditambahkan secara bertahap untuk klasifikasi tiga arah H1 pada bar dan periode uji yang sama.

\textcite{hisyam2025implementasi} menggunakan XGBoost, optimasi Bayesian, dan \emph{sliding walk-forward validation} 14 hari untuk meramalkan harga penutupan dan \emph{buyback} emas Antam harian dalam rupiah per gram. RSI, EMA, dan MACD dianalisis melalui kepentingan fitur. Sementara itu, \textcite{arif2026multivariate} menggunakan LSTM untuk meramalkan harga penutupan XAUUSD H1 tahun 2024 berdasarkan data OHLC dan lima pemisahan \emph{rolling-origin}. Kedua studi mendukung penggunaan XGBoost, indikator teknikal, objek XAUUSD, dan evaluasi temporal pada konteks masing-masing. Penelitian ini berfokus pada perubahan skor antarkelompok fitur setelah model dilatih dan disetel ulang untuk klasifikasi tiga kelas.

\subsection{Perbandingan Studi Terdekat dan Posisi Penelitian}

Tabel \ref{tab:studi-terdekat} merangkum tujuh penelitian yang paling dekat berdasarkan objek, target, fitur, model, atau cara evaluasinya. Perbedaan yang disajikan memperjelas posisi penelitian ini terhadap tujuan setiap studi.

\begingroup
\setlength{\LTleft}{\fill}
\setlength{\LTright}{\fill}
\setlength{\tabcolsep}{3pt}
\begin{longtable}{>{\raggedright\arraybackslash}p{2.2cm}>{\raggedright\arraybackslash}p{2.8cm}>{\raggedright\arraybackslash}p{3.4cm}>{\raggedright\arraybackslash}p{4.2cm}}
    \caption{Perbandingan tujuh penelitian terdekat}
    \label{tab:studi-terdekat}\\
    \hline
    Studi & Target dan frekuensi & Metode dan fitur utama & Perbedaan dengan penelitian ini \\
    \hline
    \endfirsthead
    \caption[]{Perbandingan tujuh penelitian terdekat (lanjutan)}\\
    \hline
    Studi & Target dan frekuensi & Metode dan fitur utama & Perbedaan dengan penelitian ini \\
    \hline
    \endhead
    \hline
    \multicolumn{4}{r}{Bersambung ke halaman berikutnya}\\
    \endfoot
    \hline
    \endlastfoot
    Sadorsky (2021) & Arah biner ETF GLD/SLV harian; horizon 1--20 hari & Logit, bagging, boosting, random forest; 13 indikator teknikal & Harga spot XAUUSD H1; tiga kelas; evaluasi \emph{walk-forward}. \\
    Basher dan Sadorsky (2022) & Arah emas dan Bitcoin harian & Logit, bagging, random forest; fitur teknikal dan makro; validasi deret waktu & BUY/HOLD/SELL H1; tiga kelompok fitur bertahap. \\
    Mamoudan dkk. (2023) & Sinyal \emph{FAKE}/\emph{REAL}; frekuensi tidak dinyatakan & CNN--BiGRU dan metaheuristik; MACD, MA, Ichimoku & Arah harga satu bar berikutnya; dua model dengan karakter berbeda. \\
    Taneva-Angelova dkk. (2025) & Harga emas kontinu harian; 2014--2024 & Fusi finansial, makro, sentimen; ekonometrika, ML, DL & Klasifikasi H1; macro-F1 dan balanced accuracy. \\
    Velez Flores dkk. (2026) & Harga emas kontinu bulanan; Jan. 2020--Jun. 2026 & 32 kombinasi model--fitur; validasi temporal & Tiga kelas H1; penambahan kelompok pada bar uji yang sama. \\
    Hisyam dkk. (2025) & Harga Antam harian dalam rupiah per gram & XGBoost, optimasi Bayesian, \emph{sliding walk-forward}; RSI, EMA, MACD & XAUUSD H1; selisih skor kelompok fitur. \\
    Arif dkk. (2026) & Harga penutupan XAUUSD H1 tahun 2024 & LSTM multivariat, OHLC, lima pemisahan \emph{rolling-origin} & Klasifikasi tiga kelas; penambahan teknikal dan makro; dua model. \\
\end{longtable}
\endgroup

Dalam batas studi yang ditelaah, unsur klasifikasi arah, fitur teknikal dan makroekonomi, perbandingan konfigurasi fitur, XGBoost, XAUUSD H1, serta validasi temporal telah digunakan dalam penelitian terdahulu. Celah yang dikaji terletak pada penggabungan unsur-unsur tersebut untuk klasifikasi BUY, HOLD, dan SELL XAUUSD H1. Penelitian ini akan membandingkan tiga konfigurasi fitur kumulatif menggunakan model linear dan nonlinear pada observasi serta periode uji kronologis yang sama.

Kontribusi fitur akan diukur sebagai perubahan kinerja yang bersifat kondisional terhadap urutan penambahan kelompok, model, periode, dan penyetelan. Pada analisis utama, setiap konfigurasi akan disetel secara terpisah sehingga selisih skor turut mencerminkan respons penyetelan ulang. Analisis parameter tetap pada Bab 3 akan digunakan untuk memeriksa sensitivitas hasil. Dengan demikian, hasil penelitian akan dibaca sebagai temuan prediktif dan deskriptif dalam ruang lingkup pengujian.

\section{Landasan Teori}

\subsection{XAUUSD, Bar H1, dan Waktu Keputusan}

XAUUSD menyatakan harga satu troy ounce emas dalam dolar Amerika Serikat. Kenaikan XAUUSD menunjukkan penguatan emas relatif terhadap dolar Amerika Serikat, sedangkan penurunan menunjukkan keadaan sebaliknya. Data H1 merangkum satu jam perdagangan melalui harga pembukaan $O_t$, tertinggi $H_t$, terendah $L_t$, dan penutupan $C_t$. MetaTrader 5 juga dapat menyediakan \emph{tick volume}, yaitu jumlah tick atau pembaruan kuotasi yang tercatat, yang secara konseptual berbeda dari volume transaksi pada bursa terpusat. Kolom yang akan digunakan ditetapkan berdasarkan hasil pemeriksaan data pada Bab 3.

Pada akhir bar ke-$t$, riwayat sampai $C_t$ telah tersedia untuk membentuk fitur, sedangkan $C_{t+1}$ baru diketahui setelah bar berikutnya selesai. Pemisahan waktu fitur, keputusan, dan akhir horizon target mencegah kebocoran informasi masa depan. Imbal hasil logaritmik pada bar ke-$t$ dihitung menggunakan Persamaan \ref{eq:logreturn} untuk merepresentasikan perubahan harga secara relatif \parencite{tsay2010analysis}.
\begin{equation}
r_t=\ln\left(\frac{C_t}{C_{t-1}}\right)
\label{eq:logreturn}
\end{equation}
Pada persamaan tersebut, $r_t$ adalah imbal hasil logaritmik bar ke-$t$, $C_t$ adalah harga penutupan bar ke-$t$, dan $C_{t-1}$ adalah harga penutupan bar sebelumnya. Besaran relatif lain, seperti bentuk bar $\ln(C_t/O_t)$, rentang $\ln(H_t/L_t)$, dan volatilitas historis, digunakan untuk merangkum dinamika harga tanpa bergantung langsung pada tingkat harga. Seluruh jendela perhitungan berakhir paling lambat pada waktu keputusan.

\subsection{Target Tiga Kelas dan Ambang Historis}

Target penelitian memetakan imbal hasil satu bar berikutnya ke dalam kelas BUY, HOLD, atau SELL. Ketiga istilah yang lazim dalam konteks pasar tersebut digunakan sebagai label arah yang dibentuk dari harga, bukan label bawaan penyedia data maupun instruksi transaksi. Dalam pemodelan pembelajaran mesin finansial, penggunaan ambang batas dinamis berbasis volatilitas lokal bertujuan memisahkan fluktuasi tren yang bermakna dari derau acak pasar \parencite{lopezdeprado2018advances}. Dengan ambang volatilitas $\delta_t$, aturan pembentukan kelas ditunjukkan pada Persamaan \ref{eq:labeling}.
\begin{equation}
y_t=\begin{cases}\mathrm{BUY},&r_{t+1}>\delta_t\\\mathrm{HOLD},&-\delta_t\le r_{t+1}\le\delta_t\\\mathrm{SELL},&r_{t+1}<-\delta_t\end{cases}
\label{eq:labeling}
\end{equation}
Pada persamaan tersebut, $y_t$ adalah label untuk keputusan pada bar ke-$t$, $r_{t+1}$ adalah imbal hasil bar berikutnya, dan $\delta_t$ adalah ambang yang dihitung hanya dari data historis sampai bar ke-$t$. Kelas HOLD merepresentasikan perubahan yang kecil relatif terhadap volatilitas historis. Horizon dan rumus operasional ambang dijelaskan pada Bab 3.

Seluruh konfigurasi akan menggunakan label dan bar layak yang sama. Masker observasi bersama diperlukan agar selisih skor mencerminkan perubahan fitur, bukan perubahan sampel. Ketentuan yang sama berlaku pada kedua model sehingga perbandingan dapat dilakukan secara berpasangan.

\subsection{Fitur Teknikal sebagai Transformasi Riwayat}

Indikator teknikal meringkas riwayat harga menjadi representasi yang dapat dipelajari oleh model. \emph{Relative Strength Index} (RSI) menggambarkan keseimbangan kenaikan dan penurunan harga, sedangkan \emph{Moving Average Convergence Divergence} (MACD) menggambarkan perbedaan tren eksponensial jangka pendek dan jangka panjang \parencite{wilder1978new,murphy1999technical}. Keduanya merupakan transformasi deterministik dari riwayat harga, sehingga manfaatnya dinilai dari kemampuan representasi tersebut dalam membantu klasifikasi.

Untuk periode RSI $p$, perubahan harga dinyatakan sebagai $d_t=C_t-C_{t-1}$, kenaikan sebagai $g_t=\max(d_t,0)$, dan penurunan sebagai $\ell_t=\max(-d_t,0)$. Setelah rata-rata awal dari $p$ perubahan, penghalusan Wilder dihitung sebagai berikut.
\begin{align}
\overline g_t&=\frac{(p-1)\overline g_{t-1}+g_t}{p}, &
\overline \ell_t&=\frac{(p-1)\overline \ell_{t-1}+\ell_t}{p},\\
RS_t&=\frac{\overline g_t}{\overline \ell_t}, &
RSI_t&=100-\frac{100}{1+RS_t}.
\end{align}
Simbol $p$ menyatakan panjang periode, sedangkan $\overline g_t$ dan $\overline \ell_t$ masing-masing menyatakan rata-rata kenaikan dan penurunan yang dihaluskan pada bar ke-$t$. Jika kedua rata-rata bernilai nol, RSI ditetapkan 50. Jika hanya $\overline \ell_t$ yang nol, RSI ditetapkan 100; jika hanya $\overline g_t$ yang nol, RSI ditetapkan 0.

Rata-rata bergerak eksponensial atau \emph{exponential moving average} (EMA) dengan periode $p$ menggunakan faktor penghalusan $\alpha_p=2/(p+1)$. Perhitungan EMA, MACD, garis sinyal, dan histogram ditunjukkan sebagai berikut.
\begin{align}
\operatorname{EMA}_{p,t}&=\alpha_p C_t+(1-\alpha_p)\operatorname{EMA}_{p,t-1},\\
\operatorname{MACD}_t&=\operatorname{EMA}_{12,t}-\operatorname{EMA}_{26,t},\\
\operatorname{Signal}_t&=\operatorname{EMA}_{9}(\operatorname{MACD})_t,\\
\operatorname{Histogram}_t&=\operatorname{MACD}_t-\operatorname{Signal}_t.
\end{align}
Pada persamaan tersebut, $\operatorname{EMA}_{p,t}$ adalah EMA periode $p$ pada bar ke-$t$, $\operatorname{Signal}_t$ adalah EMA sembilan periode dari MACD, dan $\operatorname{Histogram}_t$ adalah selisih MACD dengan garis sinyal. Nilai awal dan masa pemanasan memengaruhi bagian awal deret, sehingga inisialisasi, normalisasi, dan masa pemanasan akan ditetapkan pada data pelatihan tanpa menggunakan data uji.

\subsection{Fitur Makroekonomi dan Ketertiban Waktu}

Indikator makroekonomi menggambarkan kondisi ekonomi agregat. \emph{Consumer Price Index} (CPI) mengukur perubahan harga sekeranjang barang dan jasa konsumsi, sedangkan CPI inti mengecualikan pangan dan energi. Perubahan bulanan membandingkan indeks dengan bulan sebelumnya, sementara perubahan tahunan membandingkannya dengan bulan yang sama setahun sebelumnya.\footnote{Penjelasan CPI oleh BLS: \url{https://www.bls.gov/cpi/questions-and-answers.htm}.} \emph{Nonfarm Payrolls} (NFP) merujuk pada pekerjaan penggajian nonpertanian; rilis perubahan bulanannya menggambarkan perubahan jumlah pekerjaan. Tingkat pengangguran menyatakan persentase angkatan kerja yang menganggur dan berasal dari survei rumah tangga, sedangkan NFP berasal dari survei perusahaan.\footnote{Catatan teknis statistik ketenagakerjaan BLS: \url{https://www.bls.gov/news.release/empsit.tn.htm}.}

\emph{Federal Open Market Committee} (FOMC) menetapkan arah kebijakan moneter Amerika Serikat, termasuk sasaran suku bunga federal funds.\footnote{Penjelasan FOMC oleh Federal Reserve: \url{https://www.federalreserve.gov/monetarypolicy/fomc.htm}.} Keputusan suku bunga dan publikasi FOMC diperlakukan sebagai peristiwa kalender yang berbeda. Informasi inflasi, ketenagakerjaan, dan kebijakan moneter memberikan konteks bagi pengujian fitur eksternal pada emas. Arah serta besar kontribusi prediktifnya akan diukur melalui eksperimen, tanpa menetapkan bahwa kenaikan suatu indikator selalu diikuti arah harga emas tertentu.

Fitur makroekonomi berasal dari kalender ekonomi terstruktur. Salah satu bentuk utamanya adalah kejutan makro (\emph{macro surprise}), yaitu selisih antara nilai aktual dan prakiraan ($actual - forecast$). Nilai ini menggambarkan perbedaan hasil pengumuman terhadap nilai prakiraan yang tercatat pada kalender. Untuk menjaga urutan waktu kausal dan mencegah kebocoran informasi masa depan, waktu rilis yang telah dikonversi ke UTC harus lebih kecil secara ketat daripada waktu keputusan.

Dalam penelitian ini, konfigurasi K3 dievaluasi menggunakan arsip kalender MetaTrader 5 yang telah dibekukan. Kolom \texttt{forecast} diperlakukan secara operasional sebagai nilai prakiraan yang diasosiasikan vendor dengan peristiwa rilis, sedangkan kolom \texttt{actual} diperlakukan sebagai nilai aktual pada arsip tersebut. Peran semantik kedua kolom ini dibedakan dari verifikasi \emph{point-in-time} independen karena riwayat perubahan nilai tidak tersedia. Oleh sebab itu, kontribusi prediktif K3 ditafsirkan sebagai efektivitas fitur dari snapshot kalender yang digunakan, bukan sebagai pembuktian dampak kausal makroekonomi atau simulasi waktu nyata yang sepenuhnya bebas dari potensi revisi data.

Keadaan tanpa rilis, nilai yang hilang, dan informasi yang kedaluwarsa dibedakan dari kejutan bernilai nol. Kejutan nol berarti nilai aktual sama dengan prakiraan, sedangkan nilai hilang menunjukkan bahwa besarnya kejutan tidak diketahui. Umur rilis dan indikator ketiadaan digunakan untuk mempertahankan perbedaan tersebut. Nilai lama juga tidak diteruskan ketika terdapat rilis yang lebih baru pada keluarga peristiwa yang sama, sehingga informasi lama tidak dipertahankan seolah-olah masih merupakan rilis terkini.

\subsection{Konfigurasi dan Makna Kontribusi Bersyarat}

Tiga konfigurasi fitur dibentuk secara kumulatif. K1 memuat fitur harga, K2 menambahkan fitur teknikal, dan K3 menambahkan fitur makroekonomi dari kalender MetaTrader 5. Untuk metrik $S$, perubahan kinerja akibat penambahan kelompok teknikal dan makroekonomi dihitung dengan Persamaan \ref{eq:delta_teknikal} dan \ref{eq:delta_makro}.
\begin{equation}
\Delta S_{\mathrm{teknikal}}=S_2-S_1
\label{eq:delta_teknikal}
\end{equation}
\begin{equation}
\Delta S_{\mathrm{makro}}=S_3-S_2
\label{eq:delta_makro}
\end{equation}
Simbol $S_1$, $S_2$, dan $S_3$ menyatakan skor K1, K2, dan K3 pada model, metrik, periode uji, label, serta bar uji yang sama. Nilai perubahan dapat positif, nol, atau negatif.

Urutan K1--K2--K3 membuat kontribusi bersifat kondisional. Selisih $S_3-S_2$, misalnya, mengukur perubahan setelah fitur makroekonomi ditambahkan pada fitur harga dan teknikal. Pada analisis utama, setiap konfigurasi akan disetel secara terpisah sehingga selisih juga mencakup respons terhadap penyetelan ulang. Analisis parameter tetap digunakan untuk memeriksa sensitivitas terhadap pilihan tersebut. Pengukuran dilakukan pada tingkat kelompok fitur dan ditafsirkan sebagai perubahan prediktif dalam ruang lingkup penelitian.

\subsection{Regresi Logistik Multinomial}

Regresi Logistik Multinomial memodelkan probabilitas lebih dari dua kelas dan digunakan sebagai model linear \parencite{hosmer2013applied}. Untuk vektor fitur $\mathbf{x}$, probabilitas kelas dihitung dengan fungsi softmax pada Persamaan \ref{eq:softmax}.
\begin{equation}
P(y=k\mid\mathbf{x})=\frac{\exp\left(\beta_{0k}+\boldsymbol{\beta}_k^{\mathsf{T}}\mathbf{x}\right)}{\sum_{j=1}^{K}\exp\left(\beta_{0j}+\boldsymbol{\beta}_j^{\mathsf{T}}\mathbf{x}\right)}
\label{eq:softmax}
\end{equation}
Simbol $K=3$ menyatakan jumlah kelas, $\beta_{0k}$ menyatakan intersep kelas ke-$k$, dan $\boldsymbol{\beta}_k$ menyatakan vektor koefisien kelas ke-$k$. Prediksi ditentukan dari kelas dengan probabilitas terbesar.

Ketidakseimbangan kelas ditangani melalui bobot kelas pada fungsi objektif. Dengan penalti L2, tujuan pelatihan dituliskan pada Persamaan \ref{eq:logistic_objective}.
\begin{equation}
\min_{\boldsymbol{\beta}}\left[-\sum_{i=1}^{n}w_{y_i}\sum_{k=1}^{K}\mathbf{1}(y_i=k)\log P(y_i=k\mid\mathbf{x}_i)+\lambda\sum_{k=1}^{K}\lVert\boldsymbol{\beta}_k\rVert_2^2\right]
\label{eq:logistic_objective}
\end{equation}
Simbol $n$ menyatakan jumlah data pelatihan, $w_{y_i}$ adalah bobot kelas observasi ke-$i$, $\mathbf{1}(y_i=k)$ adalah fungsi indikator, dan $\lambda$ adalah kekuatan regularisasi. Pembobotan mengurangi dominasi kelas yang lebih sering muncul, sedangkan penalti L2 mengendalikan besar koefisien. Pada implementasi, hiperparameter $C$ merupakan kebalikan kekuatan regularisasi sehingga nilai $C$ yang lebih kecil menghasilkan regularisasi lebih kuat. Parameter standardisasi dipelajari hanya dari data pelatihan pada setiap pemisahan latih-validasi atau latih-uji.

\subsection{XGBoost untuk Klasifikasi Multikelas}

XGBoost membentuk prediktor aditif dari sekumpulan pohon keputusan. Vektor skor untuk observasi ke-$i$ dinyatakan pada Persamaan \ref{eq:xgboost_score}.
\begin{equation}
\widehat{\mathbf{s}}_i=\sum_{m=1}^{M}\mathbf{f}_m(\mathbf{x}_i)
\label{eq:xgboost_score}
\end{equation}
Simbol $\widehat{\mathbf{s}}_i$ menyatakan vektor skor untuk tiga kelas, $M$ adalah jumlah tahap boosting, dan $\mathbf{f}_m$ adalah keluaran kelompok pohon pada tahap ke-$m$, termasuk pengaruh laju belajar. Objektif pelatihan menggabungkan kerugian klasifikasi dan penalti kompleksitas pohon, seperti jumlah daun dan besar bobot daun \parencite{chen2016xgboost}.\footnote{Dokumentasi model XGBoost: \url{https://xgboost.readthedocs.io/en/stable/tutorials/model.html}.}

Pada klasifikasi multikelas, vektor skor diubah menjadi probabilitas dan kelas dengan probabilitas terbesar dipilih sebagai prediksi. Struktur pohon dapat menangkap ambang nonlinear dan interaksi antarfitur. Fleksibilitas ini dikendalikan melalui hiperparameter kedalaman, laju belajar, jumlah tahap boosting, subsampling, dan regularisasi yang dipilih pada validasi internal. Dalam penelitian ini, XGBoost digunakan sebagai model nonlinear berbasis pohon.

\subsection{Evaluasi Kronologis dan Pencegahan Kebocoran}

Evaluasi deret waktu mengikuti urutan penggunaan data: periode yang lebih awal digunakan untuk pelatihan dan periode berikutnya digunakan untuk pengujian. Rancangan \emph{nested walk-forward} menggunakan blok validasi internal pada data pelatihan untuk memilih hiperparameter, kemudian model dilatih ulang dan dinilai satu kali pada periode uji utama yang lebih baru. Imputasi, seleksi kolom, penskalaan, bobot kelas, dan penyetelan dipelajari kembali hanya dari bagian pelatihan.

\emph{Purge} menghapus observasi pelatihan yang horizon labelnya bertumpang tindih dengan awal validasi atau pengujian. Sementara itu, embargo merupakan jeda konservatif tambahan setelah batas pemisahan. Riwayat harga dalam jendela fitur telah tersedia pada waktu keputusan, sehingga panjang \emph{purge} ditentukan oleh horizon label. Aturan ini mencegah label pelatihan yang berakhir pada atau setelah keputusan pertama pada blok berikutnya.

Perbandingan konfigurasi dilakukan pada bar, label, periode uji, dan metrik yang sama. Prediktor kelas mayoritas dari data pelatihan digunakan sebagai diagnostik kewajaran. Skor akan dilaporkan per periode uji dan dibandingkan secara berpasangan. Karena periode uji dengan jendela mengembang saling bergantung, ringkasan antarperiode uji ditafsirkan secara deskriptif.

\subsection{Macro-F1 dan Balanced Accuracy}

Matriks kebingungan (\emph{confusion matrix}) merangkum kesesuaian kelas sebenarnya dengan kelas prediksi. Pada Tabel \ref{tab:confusion-tiga-kelas}, $n_{ij}$ adalah jumlah observasi berkelas sebenarnya $i$ yang diprediksi sebagai $j$; $B$, $H$, dan $S$ masing-masing berarti BUY, HOLD, dan SELL. Tabel ini menunjukkan bentuk perhitungan, bukan hasil eksperimen.

\begin{table}[htbp]
\centering\caption{Bentuk matriks kebingungan tiga kelas}\label{tab:confusion-tiga-kelas}
\begin{tabular}{lccc}\hline
Kelas sebenarnya & Prediksi BUY & Prediksi HOLD & Prediksi SELL\\\hline
BUY & $n_{BB}$ & $n_{BH}$ & $n_{BS}$\\
HOLD & $n_{HB}$ & $n_{HH}$ & $n_{HS}$\\
SELL & $n_{SB}$ & $n_{SH}$ & $n_{SS}$\\\hline
\end{tabular}
\end{table}

Untuk setiap kelas $c$, dua kelas lainnya diperlakukan sebagai kelas pembanding. Sebagai contoh, $TP_B=n_{BB}$, $FP_B=n_{HB}+n_{SB}$, dan $FN_B=n_{BH}+n_{BS}$. Dengan demikian, diagonal matriks berisi prediksi yang benar, sedangkan sel di luar diagonal menunjukkan jenis kesalahan klasifikasi.

Untuk kelas $c$, precision, recall, dan F1 dihitung menggunakan Persamaan \ref{eq:precision}--\ref{eq:f1}.
\begin{equation}
\operatorname{Precision}_c=\frac{TP_c}{TP_c+FP_c}
\label{eq:precision}
\end{equation}
\begin{equation}
\operatorname{Recall}_c=\frac{TP_c}{TP_c+FN_c}
\label{eq:recall}
\end{equation}
\begin{equation}
F1_c=\frac{2TP_c}{2TP_c+FP_c+FN_c}
\label{eq:f1}
\end{equation}
Simbol $TP_c$, $FP_c$, dan $FN_c$ masing-masing menyatakan jumlah positif benar, positif salah, dan negatif salah untuk kelas $c$ \parencite{sokolova2009systematic}. Macro-F1 memberi bobot yang sama kepada ketiga kelas seperti ditunjukkan pada Persamaan \ref{eq:macrof1}.
\begin{equation}
\operatorname{Macro\text{-}F1}=\frac{F1_{\mathrm{BUY}}+F1_{\mathrm{HOLD}}+F1_{\mathrm{SELL}}}{3}
\label{eq:macrof1}
\end{equation}
Jika penyebut F1 suatu kelas bernilai nol, F1 kelas tersebut ditetapkan nol dan penyebut macro-F1 tetap tiga kelas. Konvensi ini memungkinkan kinerja setiap kelas tetap terwakili ketika suatu kelas tidak pernah diprediksi.

Balanced accuracy merupakan rata-rata recall kelas dan dihitung dengan Persamaan \ref{eq:balancedacc}.
\begin{equation}
\begin{split}
\operatorname{Balanced\ Accuracy} &= \frac{1}{3} \big( \operatorname{Recall}_{\mathrm{BUY}} \\
&\quad + \operatorname{Recall}_{\mathrm{HOLD}} + \operatorname{Recall}_{\mathrm{SELL}} \big).
\end{split}
\label{eq:balancedacc}
\end{equation}
Jika suatu kelas aktual tidak muncul pada periode uji, recall kelas tersebut dan balanced accuracy tiga kelas tidak terdefinisi. Nilainya tidak diganti dengan nol. Macro-F1 tetap dihitung menggunakan konvensi penyebut tiga kelas, sehingga kedua metrik digunakan secara terpisah untuk membaca perubahan kinerja.

\chapter{METODOLOGI}

Bab ini menjelaskan protokol penelitian yang diusulkan untuk mengukur perubahan kinerja klasifikasi arah XAUUSD setelah penambahan kelompok fitur teknikal dan makroekonomi. Seluruh ambang, periode, fitur, dan hiperparameter ditetapkan sebagai pilihan desain sebelum eksperimen utama, bukan sebagai nilai optimum yang telah dibuktikan. Arsip uji kelayakan yang tersedia menjadi bukti bahwa penelitian dapat dilanjutkan, tetapi hasilnya tidak diperlakukan sebagai hasil eksperimen utama ataupun sebagai data uji prospektif yang belum pernah dilihat.

\section{Tipe Penelitian}

Penelitian ini bertipe nonimplementatif--analitik dalam bidang komputasi cerdas. Tugas klasifikasi digunakan sebagai instrumen untuk membandingkan kandungan informasi tiga konfigurasi fitur, bukan untuk membangun aplikasi transaksi, memberikan rekomendasi perdagangan, atau menjelaskan hubungan sebab-akibat ekonomi. Penelitian tidak melibatkan partisipan manusia.

\section{Strategi Penelitian}

Strategi penelitian berupa eksperimen komparatif dengan penambahan kelompok fitur secara kumulatif. Dua model dengan karakter berbeda digunakan, yaitu Regresi Logistik Multinomial sebagai pembanding linear dan XGBoost sebagai pembanding nonlinear berbasis pohon. Pemakaian kedua model tersebut mengikuti preseden klasifikasi arah aset dengan model linear dan berbasis pohon \parencite{sadorsky2021predicting,basher2022forecasting}. Kedua model diterapkan pada K1, K2, dan K3 untuk menilai apakah perubahan kinerja antarkonfigurasi konsisten pada karakter model yang berbeda. Perbedaan skor kedua model dilaporkan secara deskriptif, bukan sebagai uji linearitas data atau dasar untuk menyatakan satu model selalu lebih unggul.

Perubahan skor Konfigurasi 2 terhadap Konfigurasi 1 ($K2-K1$) dan Konfigurasi 3 terhadap Konfigurasi 2 ($K3-K2$) dihitung secara berpasangan pada bar uji, periode uji, model, dan metrik yang sama. Pada analisis utama, setiap konfigurasi memperoleh penyetelan hiperparameter sendiri menggunakan data latih. Dengan demikian, selisih skor mengukur perubahan bersyarat akibat penambahan kelompok fitur beserta penyetelan ulang model, bukan pengaruh murni atau kausal fitur tersebut. Analisis sensitivitas dengan hiperparameter tetap digunakan untuk membantu memisahkan kedua unsur itu.

Alur penelitian yang direncanakan ditunjukkan pada Gambar \ref{fig:alur-penelitian}. Protokol dibekukan sebelum eksperimen baru dijalankan dan sebelum skor periode uji diperiksa.

\begin{figure}[htbp]
    \centering
    \begin{tikzpicture}[
        node distance=3.5mm,
        every node/.style={font=\small,align=center},
        process/.style={rectangle,draw=black,rounded corners=1.5pt,
            minimum width=10.8cm,minimum height=0.83cm,text width=10.2cm,
            inner xsep=4pt,inner ysep=3pt},
        flow/.style={-{Latex[length=2.2mm]},thin}
    ]
        \node[process] (input) {Mulai: data harga dan kalender ekonomi sekunder};
        \node[process,below=of input] (audit) {Audit kualitas harga dan kalender\\serta verifikasi integritas dan waktu rilis};
        \node[process,below=of audit] (feature) {Pelabelan kronologis dan pembentukan fitur\\harga, teknikal, dan makroekonomi};
        \node[process,below=of feature] (inner) {Validasi internal: praproses, penyetelan,\\dan pelatihan ulang pada data latih luar};
        \node[process,below=of inner] (outer) {Prediksi pada periode uji utama dan penghitungan\\macro-F1, balanced accuracy, serta matriks kebingungan};
        \node[process,below=of outer] (compare) {Perbandingan berpasangan $K2-K1$ dan $K3-K2$\\serta analisis sensitivitas};
        \draw[flow] (input) -- (audit);
        \draw[flow] (audit) -- (feature);
        \draw[flow] (feature) -- (inner);
        \draw[flow] (inner) -- (outer);
        \draw[flow] (outer) -- (compare);
    \end{tikzpicture}
    \caption{Alur penelitian yang diusulkan}
    \label{fig:alur-penelitian}
\end{figure}

\section{Sumber dan Teknik Pengumpulan Data}

\subsection{Sumber dan Cakupan Data}

Penelitian menggunakan data sekunder yang direkam oleh vendor. Data harga utama berupa bar \emph{open}, \emph{high}, \emph{low}, dan \emph{close} (OHLC) XAUUSD H1 dari server MetaQuotes-Demo. Arsip kelayakan yang telah diverifikasi mencakup 38.799 bar dari Januari 2020 sampai Juli 2026 atau sekitar 6,578 tahun. Waktu bar yang dikembalikan antarmuka Python MetaTrader 5 diperlakukan sebagai UTC sesuai dokumentasi resminya.\footnote{Dokumentasi resmi \emph{copy rates range}: \url{https://www.mql5.com/en/docs/python_metatrader5/mt5copyratesrange_py}.} Jumlah observasi eksperimen akhir ditentukan setelah seluruh aturan kelayakan, jeda, pemanasan indikator, dan pelabelan diterapkan.

Data makroekonomi berasal dari kalender ekonomi internal MetaTrader 5. Snapshot arsip beku ini berisi 706 kejadian USD yang terverifikasi tanpa duplikasi pengenal nilai, dan keluarga kejadian seperti NFP dan tingkat pengangguran memiliki data aktual dan prakiraan utuh yang secara struktural menutupi periode data latih dan setiap periode uji utama. Dokumentasi kalender menyatakan bahwa waktu kalender menggunakan waktu server perdagangan.\footnote{Dokumentasi resmi \emph{CalendarValueHistory}: \url{https://www.mql5.com/en/docs/calendar/calendarvaluehistory}.} Namun, apabila penanda waktu mentah historis pada snapshot beku spesifik ini dilabeli ulang sebagai UTC, seluruh sepuluh jangkar waktu rilis resmi telah terverifikasi cocok secara tepat; oleh sebab itu, rentang waktu arsip kalender ini diperlakukan sebagai UTC dan dikontrol ulang menggunakan pemeriksa hash SHA-256 dan penanda waktu pada Subbab \ref{subsec:audit-kalender}.

Keluarga peristiwa dibatasi pada CPI umum bulanan, CPI umum tahunan, CPI inti bulanan, CPI inti tahunan, \emph{Nonfarm Payrolls} (NFP), tingkat pengangguran, keputusan suku bunga, dan peristiwa FOMC. Pemilihan ini didukung oleh temuan bahwa kejutan makroekonomi berkaitan dengan pembentukan harga emas intrahari \parencite{sobti2021how}.

\subsection{Teknik Pengumpulan Data}

Bar harga diperoleh secara baca-saja melalui antarmuka Python MetaTrader 5. Ekstraksi arsip awal dilakukan pada Windows 11 tanpa pemanggilan API pesanan atau perdagangan. Hasil ekstraksi disimpan sebagai CSV dan Parquet bersama manifes yang mencatat sumber, simbol, rentang waktu, serta checksum berkas. Windows dan terminal MT5 hanya diperlukan kembali apabila ekstraktor harga Python digunakan untuk mengambil sumber atau rentang harga baru; penyuntingan proposal dan analisis luring atas arsip yang tersedia tidak memerlukannya.

Kalender ekonomi diperoleh secara baca-saja melalui eksportir MQL5 yang telah berfungsi pada MT5 dalam lingkungan Wine di macOS. CSV hasil ekspor mempertahankan kolom waktu mentah, nilai aktual, dan nilai prakiraan agar provenance serta pemeriksaan waktu dapat dilakukan. Pengambilan ini tidak menggunakan kredensial di dalam kode, tidak mengaktifkan transaksi otomatis, dan tidak memanggil API pesanan.

Data harga dan kalender digunakan secara privat dan luring. Berkas mentah dipertahankan terpisah dari data turunan; keluaran pemrosesan tidak menimpa sumber. Sebelum digunakan, identitas sumber, rentang waktu, skema, jumlah bar atau kejadian, checksum, duplikasi, urutan waktu, dan kelengkapan kolom diperiksa serta dicatat dalam manifes audit.

\section{Lingkungan dan Peralatan Penelitian}

Penelitian tidak memerlukan lokasi fisik atau laboratorium tertentu karena penyiapan dan analisis data dilakukan secara luring pada komputer lokal. Komputasi dirancang untuk CPU dan tidak mensyaratkan GPU. Lingkungan persiapan yang tersedia telah diukur pada 12 September 2026 sebagaimana Tabel \ref{tab:perangkat-keras}; pencatatan ini tidak menyatakan bahwa eksperimen utama telah dilaksanakan.

\begin{table}[htbp]
    \centering
    \caption{Perangkat keras lingkungan persiapan penelitian}
    \label{tab:perangkat-keras}
    \begin{tabular}{ll}
        \hline
        Komponen & Spesifikasi terukur \\
        \hline
        Prosesor & Apple M5 (arsitektur arm64) \\
        Memori & 24 GiB RAM \\
        Akselerator & Tidak diperlukan; analisis berbasis CPU \\
        \hline
    \end{tabular}
\end{table}

Perangkat lunak lingkungan persiapan diringkas pada Tabel \ref{tab:perangkat-lunak}. Versi sistem operasi, Python, pustaka, dan Git pada pelaksanaan akhir akan dicatat kembali bersama artefak eksperimen. Lingkungan sumber historis arsip harga tetap dicatat secara terpisah, yaitu Windows 11 dengan Python 3.12.10, NumPy 2.5.1, pandas 3.0.5, scikit-learn 1.9.0, dan XGBoost 2.1.4; pencatatan lingkungan macOS tidak menggantikan provenance tersebut.

\begin{table}[htbp]
    \centering
    \caption{Perangkat lunak lingkungan persiapan penelitian}
    \label{tab:perangkat-lunak}
    \begin{tabular}{ll}
        \hline
        Komponen & Versi terukur pada 12 September 2026 \\
        \hline
        Sistem operasi & macOS 26.6.2 \\
        Bahasa pemrograman & Python 3.14.5 pada lingkungan virtual proyek \\
        Pustaka numerik dan data & NumPy 2.5.1; pandas 3.0.5; PyArrow 23.0.1 \\
        Pustaka pemodelan & scikit-learn 1.9.0; XGBoost 2.1.4 \\
        \hline
    \end{tabular}
\end{table}

\section{Unit Observasi, Waktu Keputusan, dan Target}

Penanda waktu pada data harga menyatakan waktu pembukaan bar. Misalkan $o_t$ adalah waktu pembukaan bar ke-$t$, sedangkan $O_t$, $H_t$, $L_t$, dan $C_t$ berturut-turut adalah harga pembukaan, tertinggi, terendah, dan penutupannya. Keputusan untuk observasi ke-$t$ dibuat setelah bar tersebut selesai, yaitu
\begin{equation}
\tau_t=o_t+1\text{ jam}.
\label{eq:waktu-keputusan}
\end{equation}
Semua fitur harga dan teknikal harus telah tersedia paling lambat pada $\tau_t$. Peristiwa makro hanya digunakan jika waktu rilis terverifikasinya lebih kecil secara ketat daripada $\tau_t$; rilis yang tepat sama dengan $\tau_t$ ditunda ke keputusan berikutnya.

Target merupakan imbal hasil penutupan satu bar berikutnya,
\begin{equation}
q_{t+1}=\ln\left(\frac{C_{t+1}}{C_t}\right).
\label{eq:target-satu-jam}
\end{equation}
Data harga yang diperoleh tidak memuat label BUY, HOLD, atau SELL. Peneliti membentuk label tersebut secara deterministik dari $C_t$ dan $C_{t+1}$: BUY untuk kenaikan di atas ambang positif, HOLD untuk perubahan di antara kedua ambang, dan SELL untuk penurunan di bawah ambang negatif. Ketiganya menyatakan kelas arah, bukan perintah transaksi. Rumus ambang dan syarat kelayakan bar dijelaskan berikut ini.

Observasi hanya layak diberi target apabila bar berikutnya dibuka tepat pada $o_{t+1}=o_t+1$ jam dan target berakhir pada $\tau_{t+1}=o_t+2$ jam. Dengan demikian, label tidak menjembatani akhir pekan, penutupan pasar, atau celah data. Pada setiap periode uji, $\tau_{t+1}$ juga harus berada di dalam interval uji yang sama.

Ambang kelas utama ditetapkan sebesar setengah volatilitas historis 24 bar. Dengan $r_t$ sebagaimana Persamaan \ref{eq:logreturn}, rata-rata dan simpangan baku sampel dihitung sebagai
\begin{align}
\bar r_{t,24} &= \frac{1}{24}\sum_{i=0}^{23}r_{t-i},\\
\sigma_{t,24} &= \sqrt{\frac{1}{23}\sum_{i=0}^{23}\left(r_{t-i}-\bar r_{t,24}\right)^2},\\
\delta_t &= 0{,}5\sigma_{t,24}.
\label{eq:ambang-volatilitas}
\end{align}
Nilai $r_t$ sampai $r_{t-23}$ telah diketahui pada $\tau_t$. Kelas kemudian ditentukan dengan mengganti $r_{t+1}$ pada Persamaan \ref{eq:labeling} menjadi $q_{t+1}$. Observasi dengan $\sigma_{t,24}$ nol, tidak hingga, atau tidak dapat dihitung dikeluarkan. Tidak ada imputasi pada OHLC maupun label. Pengali 0,5 tidak disetel menurut konfigurasi atau periode uji. Pengali 0,25 dan 1,0 diterapkan sebagai dua analisis sensitivitas terpisah pada seluruh konfigurasi, tanpa memilih pengali berdasarkan skor terbaik.

\section{Audit dan Penyiapan Data}

\subsection{Audit Data Harga}

Data harga diperiksa terhadap duplikasi waktu, urutan waktu, nilai yang hilang, harga nonpositif, dan pelanggaran hubungan $H_t\geq\max(O_t,C_t)$ serta $L_t\leq\min(O_t,C_t)$. Audit celah menggunakan jadwal perdagangan yang berlaku pada sumber data. Penutupan pasar yang dapat dijelaskan dicatat sebagai penutupan wajar dan jendela fitur dilanjutkan dengan bar perdagangan sah sebelumnya. Dengan aturan ini, volatilitas fitur merangkum return antarbar perdagangan, termasuk perubahan melintasi penutupan wajar, sedangkan horizon target tetap tepat satu jam.

Celah yang tidak dapat dijelaskan memutus keadaan indikator sehingga perhitungan bergulir, Wilder, dan EMA dimulai kembali setelah celah. Bar OHLC tidak diinterpolasi. Aturan kelayakan ditentukan berdasarkan kualitas data dan ketersediaan informasi, bukan skor model.

\subsection{Audit Data Kalender}
\label{subsec:audit-kalender}

Sebagai kontrol reproduksibilitas untuk snapshot spesifik ini, sebelum data kalender digabungkan, berkas harus cocok dengan hash SHA-256 arsip beku yang ditetapkan pada konfigurasi. Kolom UTC baru dibuat tanpa menimpa waktu mentah. Hasil konversinya dipastikan cocok tepat dengan sepuluh jangkar rilis resmi yang mencakup periode \emph{daylight saving time} dan non-DST. Semantik waktu dan kontrol jangkar rilis historis ini diberlakukan eksklusif untuk snapshot ini dan tidak digeneralisasikan pada broker, \emph{build}, atau snapshot kalender lainnya.

Penelitian ini menggunakan asumsi operasional bahwa kolom \emph{forecast} pada rekaman MT5 merupakan nilai prakiraan yang diasosiasikan vendor dengan kejadian sebelum nilai aktual diketahui. Asumsi ini konsisten dengan dokumentasi kalender MQL5 yang membedakan nilai prakiraan dari nilai aktual yang tersedia setelah publikasi.\footnote{Dokumentasi resmi konsep kalender MQL5: \url{https://www.mql5.com/en/book/advanced/calendar/calendar_overview}.} Riwayat perubahan independen tidak tersedia sehingga snapshot ini tidak diklaim sebagai arsip \emph{point-in-time}. Nilai \emph{previous} tidak digunakan. Untuk setiap keluarga rilis, kejadian terbaru menggantikan keadaan kejadian sebelumnya. Apabila angka pada kolom nilai kejadian kosong, nilai numeriknya tidak direkayasa menjadi nol melainkan dibiarkan tidak tersedia.

\subsection{Prapemrosesan dan Ketertelusuran}

Nilai numerik yang tidak diketahui tidak diubah menjadi nol. Pada setiap blok validasi internal dan periode uji utama, median dihitung hanya dari data latih, kemudian digunakan untuk mengisi nilai hilang pada data latih dan uji bersama indikator biner ketiadaan. Kolom yang seluruhnya hilang pada data latih dibuang dan dicatat. Daftar kolom ditentukan hanya dari data latih dan dipertahankan pada data uji. Ketersediaan angka pada data uji dilaporkan sebagai cakupan evaluasi tanpa digunakan untuk memilih ulang kolom atau menyetel model.

Arsip kelayakan telah diperiksa ketika rancangan disusun. Oleh sebab itu, periode 2020--Juli 2026 tidak diklaim sebagai \emph{holdout} prospektif yang sama sekali belum pernah terlihat. Perlindungan terhadap penyesuaian setelah melihat data dilakukan dengan membekukan definisi pada bab ini sebelum eksperimen baru, menyimpan seluruh kegagalan, dan tidak mengubah protokol berdasarkan skor periode uji.

\section{Pembentukan Fitur}

\subsection{Fitur Harga}

Kelompok harga terdiri atas sembilan fitur: enam kelambanan imbal hasil $r_t,r_{t-1},\ldots,r_{t-5}$, bentuk bar $\ln(C_t/O_t)$, rentang $\ln(H_t/L_t)$, dan $\sigma_{t,24}$. Seluruhnya telah diketahui pada $\tau_t$. Nilai OHLC mentah tidak digunakan sebagai pengganti fitur yang hilang dan tidak diimputasi.

\subsection{Fitur Teknikal}

Kelompok teknikal terdiri atas empat fitur: RSI 14 periode serta garis, sinyal, dan histogram MACD $(12,26,9)$. Pemilihan RSI dan MACD mengikuti penggunaannya sebagai indikator momentum dan tren \parencite{wilder1978new,murphy1999technical}. Tidak ada SMA opsional atau pencarian panjang jendela pada data uji.

Untuk RSI, perubahan penutupan didefinisikan sebagai $d_t=C_t-C_{t-1}$, dengan kenaikan $g_t=\max(d_t,0)$ dan penurunan $\ell_t=\max(-d_t,0)$. Rata-rata kenaikan dan penurunan awal adalah rata-rata aritmetika dari 14 perubahan pertama. Nilai berikutnya diperbarui dengan penghalusan Wilder berparameter $\alpha=1/14$. RSI dihitung sebagai $100-100/(1+\overline g_t/\overline\ell_t)$. Jika kedua rata-rata nol, RSI ditetapkan 50; jika hanya rata-rata penurunan yang nol, RSI ditetapkan 100. Jika hanya rata-rata kenaikan yang nol, rumus menghasilkan RSI 0.

EMA periode $p$ diinisialisasi dengan rata-rata $p$ penutupan pertama dan diperbarui menggunakan $\alpha_p=2/(p+1)$. Garis MACD adalah $\operatorname{EMA}_{12}(C_t)-\operatorname{EMA}_{26}(C_t)$. EMA sinyal periode 9 diinisialisasi dengan rata-rata sembilan nilai MACD pertama yang valid, dan histogram adalah selisih garis MACD dengan sinyal. Ketiga fitur MACD dibagi dengan $C_t$ agar skalanya relatif terhadap harga. Setelah awal data atau setiap celah yang tidak dapat dijelaskan, observasi baru digunakan setelah sedikitnya 50 bar perdagangan valid untuk pemanasan seluruh indikator.

\subsection{Fitur Makroekonomi}

Empat tipe CPI dipertahankan sebagai keluarga terpisah, demikian pula NFP dan tingkat pengangguran. Untuk setiap keluarga berwujud numerik tersebut, nilai kejutan rilis dihitung sebagai
\begin{equation}
u_{e}=A_{e}-F_{e},
\label{eq:kejutan-makro}
\end{equation}
dengan $A_e$ sebagai kolom \emph{actual} dan $F_e$ sebagai kolom \emph{forecast} kejadian $e$ pada snapshot MT5 beku. Kejutan tidak dihitung untuk keputusan suku bunga dan peristiwa FOMC karena kolom prakiraan pada arsip tersebut tidak tersedia. Kedua tipe peristiwa tersebut hanya menyumbangkan penanda kejadian pascarilis tanpa membentuk angka kejutan rekaan.

Untuk tiap keluarga disusun nilai kejutan terbaru, umur rilis, penanda kedaluwarsa, serta indikator bahwa keputusan berada dalam 1, 6, dan 24 jam setelah rilis. Umur dibatasi pada 168 jam untuk kebutuhan representasi; pada umur $\geq168$ jam, nilai kejutan menjadi hilang dan penanda kedaluwarsa diaktifkan. Angka dari satuan berbeda tetap berada dalam kolom terpisah dan dinormalisasi hanya berdasarkan data latih. Jika beberapa pengumuman memiliki waktu sama, urutan pemrosesan ditentukan secara tetap berdasarkan keluarga dan pengenal kejadian, kemudian seluruh informasi yang sah tersedia sebelum $\tau_t$ direpresentasikan tanpa bergantung pada urutan berkas. Jadwal peristiwa masa depan tidak digunakan dalam konfigurasi utama; seluruh penanda kalender dibentuk setelah kejadian berlangsung.

\section{Konfigurasi Eksperimen}

Tiga konfigurasi pada Tabel \ref{tab:konfigurasi-fitur} diberikan kepada kedua model. K1, K2, dan K3 dilaksanakan pada setiap periode uji utama yang memenuhi persyaratan evaluasi tiga kelas. Masker observasi akhir dibentuk sekali sebelum pelatihan berdasarkan aturan data, waktu, label, dan pemanasan. Untuk menjamin perbandingan yang setara, bar uji yang digunakan harus sama untuk seluruh konfigurasi dan kedua model; K1 atau K2 tidak memperoleh bar tambahan hanya karena tidak memerlukan fitur makro.

\begin{table}[htbp]
    \centering
    \caption{Konfigurasi fitur eksperimen}
    \label{tab:konfigurasi-fitur}
    \begin{tabular}{cll}
        \hline
        Kode & Kelompok fitur & Isi utama \\
        \hline
        K1 & Harga & 9 fitur harga \\
        K2 & Harga + teknikal & K1 + RSI dan 3 fitur MACD \\
        K3 & Harga + teknikal + makro & K2 + fitur snapshot kalender MT5 \\
        \hline
    \end{tabular}
\end{table}

Prediktor konstan yang selalu memilih kelas terbanyak pada data latih dihitung sebagai diagnostik kewajaran, bukan sebagai model ketiga. Jika jumlah kelas sama, pemilihan mengikuti urutan tetap BUY, HOLD, lalu SELL. Diagnostik ini juga menggunakan bar uji yang sama.

\section{Validasi Kronologis dan Penyetelan Model}

\subsection{Validasi Luar dan Dalam}

Evaluasi utama menggunakan \emph{expanding-window walk-forward}. Interval dinyatakan setengah terbuka dalam UTC, yaitu batas awal termasuk dan batas akhir tidak termasuk. Data latih awal adalah $[1$ Januari 2020, $1$ Januari 2023$)$. Tujuh periode uji utama dicantumkan pada Tabel \ref{tab:periode-utama}. Penyebutan semester merupakan nama blok kalender; blok terakhir secara khusus mencakup Januari--Juli 2026. Setelah suatu periode selesai diuji, observasi periode tersebut dapat masuk ke data latih untuk periode berikutnya, sesuai penggunaan model secara kronologis.

\begin{table}[htbp]
    \centering
    \caption{Periode uji kronologis data utama}
    \label{tab:periode-utama}
    \begin{tabular}{cll}
        \hline
        Periode uji & Awal (termasuk) & Akhir (tidak termasuk) \\
        \hline
        1 & 1 Januari 2023 & 1 Juli 2023 \\
        2 & 1 Juli 2023 & 1 Januari 2024 \\
        3 & 1 Januari 2024 & 1 Juli 2024 \\
        4 & 1 Juli 2024 & 1 Januari 2025 \\
        5 & 1 Januari 2025 & 1 Juli 2025 \\
        6 & 1 Juli 2025 & 1 Januari 2026 \\
        7 & 1 Januari 2026 & 1 Agustus 2026 \\
        \hline
    \end{tabular}
\end{table}

Pada setiap pelatihan luar, validasi internal menggunakan 18 bulan terakhir dari data latih sebagai tiga blok validasi enam bulanan. Seluruh riwayat sebelum awal tiap blok menjadi data latih internal yang mengembang. Untuk periode uji utama pertama, blok validasinya adalah Juli--Desember 2021, Januari--Juni 2022, dan Juli--Desember 2022. Pada periode uji berikutnya, jendela 18 bulan bergerak maju mengikuti akhir data latih luar.

\subsection{Purge dan Prapemrosesan pada Setiap Pemisahan}

Sebelum setiap blok validasi atau uji, seluruh observasi latih dengan akhir label $\tau_{t+1}$ yang sama dengan atau melewati waktu keputusan pertama pada blok berikutnya dibuang. Aturan ini menghasilkan \emph{purge} sedikitnya satu bar H1 dan diterapkan pada pemisahan internal maupun luar. Pemanasan fitur dapat membaca bar lampau sebelum awal blok karena bar tersebut telah tersedia; embargo sepanjang jendela fitur tidak ditambahkan. Setiap label uji tetap wajib berakhir di dalam interval uji.

Seluruh operasi yang mempelajari parameter, yaitu imputasi median, indikator kolom hilang, pembuangan kolom yang seluruhnya hilang, penskalaan, normalisasi satuan, bobot kelas, dan pemilihan hiperparameter, diestimasi ulang hanya dari data latih pada setiap blok validasi internal dan periode uji utama. Data latih harus memuat ketiga kelas. Jika salah satu kelas tidak ada, blok atau periode tersebut dinyatakan tidak memenuhi persyaratan evaluasi tiga kelas dan tidak diubah menjadi masalah dua kelas. Kondisi tersebut dicatat sebagai kegagalan evaluasi.

\subsection{Model dan Kisi Hiperparameter}

Regresi Logistik Multinomial menggunakan penalti L2, \emph{solver} \texttt{lbfgs}, dan maksimum 2.000 iterasi. Kandidat $C$ adalah $\{0{,}01;0{,}1;1;10\}$. Fitur distandardisasi dari data latih. Bobot setiap kelas $c$ dihitung sebagai
\begin{equation}
w_c=\frac{n}{3n_c},
\label{eq:bobot-kelas}
\end{equation}
dengan $n$ jumlah observasi latih dan $n_c$ jumlah observasi kelas $c$ pada data latih.

XGBoost menggunakan objektif \texttt{multi:softprob} untuk tiga kelas. Delapan kandidat pada kisi pencarian merupakan hasil kali Kartesius parameter pada Tabel \ref{tab:kisi-xgboost}. Nilai \texttt{n\_estimators} menyatakan jumlah iterasi \emph{boosting}, bukan jumlah literal seluruh pohon pada model multikelas.

\begin{table}[htbp]
    \centering
    \caption{Kisi hiperparameter XGBoost}
    \label{tab:kisi-xgboost}
    \begin{tabular}{lll}
        \hline
        Parameter & Kandidat & Keterangan \\
        \hline
        \texttt{max\_depth} & $\{3,5\}$ & Kedalaman maksimum \\
        \texttt{learning\_rate} & $\{0{,}05;0{,}1\}$ & Laju pembelajaran \\
        \texttt{n\_estimators} & $\{100,300\}$ & Iterasi \emph{boosting} \\
        \hline
    \end{tabular}
\end{table}

Setiap kandidat XGBoost menggunakan proporsi pengambilan observasi (\texttt{subsample}) sebesar 0,8 dan proporsi fitur per pohon (\texttt{colsample\_bytree}) sebesar 0,8. Regularisasi L2 ditetapkan pada $\lambda=1$ dengan benih acak 42. Bobot observasi mengikuti Persamaan \ref{eq:bobot-kelas}. \emph{Early stopping} tidak digunakan. Kebijakan nilai hilang yang sama diterapkan melalui alur prapemrosesan Regresi Logistik dan XGBoost; kemampuan bawaan XGBoost untuk menerima nilai hilang tidak digunakan untuk mengubah definisi sampel.

\subsection{Penyetelan Hiperparameter dan Analisis Sensitivitas}

Kisi yang sama digunakan untuk setiap konfigurasi dalam model yang sama. Kandidat dipilih berdasarkan rata-rata macro-F1 pada tiga blok internal, kemudian model dipasang ulang pada seluruh data latih luar. Jika skor sama, Regresi Logistik memilih $C$ yang lebih kecil. XGBoost berturut-turut memilih kedalaman lebih kecil, jumlah iterasi \emph{boosting} lebih sedikit, lalu \emph{learning rate} lebih kecil. Kandidat yang gagal atau tidak konvergen pada salah satu blok validasi internal tidak memenuhi syarat pemilihan; rata-rata skor tidak dihitung hanya dari blok yang tersisa. Jika seluruh kandidat gagal, model pada periode uji tersebut dicatat sebagai kegagalan evaluasi beserta penyebabnya. Seluruh kegagalan tetap dicatat.

Perbandingan utama membolehkan hiperparameter berbeda antarkonfigurasi sehingga hasilnya ditafsirkan sebagai fitur ditambah penyetelan ulang. Dalam analisis sensitivitas parameter tetap, hiperparameter K1 yang dipilih pada validasi internal digunakan kembali untuk K2 dan K3 pada model dan periode uji utama yang sama. Analisis ini dilaporkan terpisah dan tidak menggantikan analisis utama.

\section{Metrik dan Analisis Hasil}

Macro-F1 menjadi metrik utama dan balanced accuracy menjadi metrik pendamping sebagaimana didefinisikan pada Bab 2. Matriks kebingungan selalu menggunakan urutan tiga label BUY, HOLD, dan SELL. Jika penyebut F1 suatu kelas nol, F1 kelas tersebut didefinisikan nol sehingga penyebut rata-rata macro-F1 tetap tiga kelas. Jika satu kelas benar tidak muncul pada data uji, recall kelas tersebut tidak terdefinisi; akibatnya balanced accuracy pada periode uji tersebut dinyatakan tidak terdefinisi dan tidak disajikan sebagai balanced accuracy tiga kelas yang sebanding. Kondisi ini, beserta jumlah kelas dan jumlah observasi, dilaporkan secara eksplisit.

Untuk setiap model, metrik, dan periode uji yang sah, perubahan skor dihitung dengan Persamaan \ref{eq:delta_teknikal} dan \ref{eq:delta_makro}. Ringkasan perubahan berpasangan meliputi rata-rata tanpa pembobotan jumlah bar, median, simpangan baku sampel, serta jumlah periode uji dengan perubahan positif. Ringkasan dihitung hanya untuk pasangan skor yang terdefinisi dan jumlah pasangan turut dilaporkan; simpangan baku tidak dihitung jika tersedia kurang dari dua pasangan. Penyebut jumlah periode uji selalu ditampilkan, termasuk periode uji yang tidak memenuhi persyaratan evaluasi, balanced accuracy yang tidak terdefinisi, dan kegagalan kandidat. Periode uji yang gagal tetap dicantumkan dalam pelaporan. Karena hanya terdapat tujuh periode uji utama yang saling bergantung akibat jendela latih mengembang, analisis bersifat deskriptif dan tidak menggunakan uji hipotesis atau nilai-$p$ naif yang mengasumsikan independensi periode uji.

Perubahan positif, nol, maupun negatif sama-sama menjawab pertanyaan penelitian. Konsistensi dinilai berdasarkan arah perubahan pada periode uji yang tersedia dan pada kedua model, tanpa mengubah kriteria kelayakan berdasarkan hasil skor. Hasil tidak ditafsirkan sebagai keuntungan transaksi, dampak kausal ekonomi, atau kontribusi individual setiap fitur.

\section{Reproduksibilitas dan Ancaman Validitas}

Setiap eksekusi menyimpan hash data masukan dan keluaran, versi perangkat lunak, benih acak, batas waktu setiap periode uji, aturan eksklusi, daftar kolom yang dibuang, parameter terpilih, status konvergensi, prediksi per bar, label benar, dan metrik. Data mentah, data turunan, dan artefak hasil ditempatkan secara terpisah sehingga proses eksperimen tidak menimpa data sumber. Catatan tersebut memungkinkan hasil ditelusuri tanpa menganggap bahwa kesamaan kode menjamin kesamaan data vendor.

Ancaman validitas utama berasal dari ketiadaan snapshot \emph{point-in-time} independen yang merekam perubahan nilai kalender. Nilai \emph{actual} dan \emph{forecast} historis yang diperoleh secara retrospektif dari MT5 mungkin mencerminkan koreksi setelah publikasi. Oleh sebab itu, hasil K3 berlaku untuk snapshot kalender beku yang digunakan dan tidak disetarakan dengan simulasi waktu nyata yang sepenuhnya \emph{point-in-time}. Ancaman lain meliputi jarangnya peristiwa makroekonomi, perubahan rezim harga, karakteristik data pada server MetaQuotes-Demo, serta ketergantungan antarperiode uji. Kesimpulan dibatasi pada XAUUSD H1 dari MetaQuotes-Demo, rentang waktu, definisi label, dua model, dan konfigurasi fitur yang dinyatakan pada penelitian ini.

\section{Jadwal Penelitian}

Jadwal penelitian disusun agar Penilaian Kemajuan 2 (P2) selesai sebelum Desember 2026. Waktu setelah P2 dialokasikan untuk menindaklanjuti masukan Pembimbing, menyelesaikan naskah dan artikel ilmiah, memenuhi persyaratan administratif, serta mempersiapkan Seminar Hasil. Eksperimen K3 dan analisis $K3-K2$ ditetapkan sebagai luaran penelitian. Rencana pelaksanaan penelitian ditunjukkan pada Tabel \ref{tab:jadwal-penelitian}.

\begin{table}[htbp]
    \centering
    \caption{Jadwal pelaksanaan penelitian}
    \label{tab:jadwal-penelitian}
    \small
    \setlength{\tabcolsep}{3pt}
    \begin{tabular}{>{\raggedright\arraybackslash}p{4.0cm}>{\raggedright\arraybackslash}p{3.0cm}>{\raggedright\arraybackslash}p{6.1cm}}
        \hline
        Kegiatan & Periode & Target \\
        \hline
        P0 dan pembekuan protokol & 21 September 2026 & Protokol serta target P1--P2 disepakati \\
        Implementasi pipeline & 22 September--4 Oktober 2026 & Pipeline validasi kronologis teruji \\
        Eksperimen K1--K2 & 5--18 Oktober 2026 & Prediksi, metrik, dan analisis sensitivitas \\
        Analisis awal dan P1 & 19--25 Oktober 2026 & Capaian sedikitnya 50\% dan ringkasan $K2-K1$ \\
        Eksperimen K3 & 26 Oktober--8 November 2026 & Prediksi, metrik, dan analisis $K3-K2$ \\
        Analisis dan penulisan & 9--22 November 2026 & Hasil, pembahasan, dan draf naskah lengkap \\
        Finalisasi dan P2 & 23--27 November 2026 & Capaian sedikitnya 80\% dan berita acara P2 \\
        Persiapan Seminar Hasil & 28 November--11 Desember 2026 & Naskah final, artikel, berkas, dan presentasi \\
        \hline
    \end{tabular}
\end{table}

\chapter{HASIL}
\label{bab:hasil}

Bab ini menyajikan hasil eksperimen klasifikasi arah XAUUSD H1 sesuai protokol Bab 3. Hasil utama menggunakan data MetaQuotes-Demo, kebijakan celah primer, dan ambang $0{,}5\sigma_{t,24}$. Hasil untuk ambang $0{,}25\sigma_{t,24}$ dan $1{,}0\sigma_{t,24}$, kebijakan \emph{holiday-reset}, serta sumber UltimaMarkets-Demo disajikan sebagai analisis sensitivitas terpisah. Angka yang dilaporkan berasal dari eksperimen Fase 4 pada 21--22 September 2026; skor uji kelayakan terdahulu tidak digunakan sebagai hasil penelitian ini.

Kinerja dilaporkan pada tujuh periode uji kronologis. Rerata skor dan perubahan antarkonfigurasi dihitung tanpa pembobotan jumlah bar, sehingga setiap periode uji mempunyai bobot sama. Penyajian menggunakan enam tempat desimal; seluruh perhitungan dilakukan dari angka sebelum pembulatan. Penafsiran hasil terhadap rumusan masalah dan keterbatasan penelitian disampaikan pada Bab 5.

\section{Hasil Audit Data dan Pembentukan Sampel}

\subsection{Cakupan Data Harga dan Celah}

Arsip harga utama berisi 38.799 bar dengan waktu pembukaan pertama 2 Januari 2020 pukul 06.00 UTC dan terakhir 31 Juli 2026 pukul 22.00 UTC. Arsip UltimaMarkets-Demo semula berisi 23.181 bar. Setelah batas awal 1 September 2022 diterapkan untuk mengecualikan lubang riwayat awal, tersisa 23.135 bar; waktu pembukaan pertama yang tersedia adalah 1 September 2022 pukul 05.00 UTC dan terakhir 31 Juli 2026 pukul 23.00 UTC. Kedua sumber dianalisis secara terpisah. Audit arsip tidak menemukan duplikasi waktu, nilai OHLC yang hilang, harga nonpositif, ataupun pelanggaran hubungan OHLC.

Tabel \ref{tab:hasil-audit-celah} merangkum klasifikasi interval antarbar yang melebihi satu jam. Penutupan terjadwal mencakup akhir pekan dan jeda sesi. Penutupan yang selaras dengan kalender bursa merupakan keselarasan antara celah teramati dan kalender bursa, dengan keterbatasan bukti historis sesi CFD broker. Kebijakan primer melanjutkan keadaan indikator setelah dua jenis penutupan tersebut. Celah tidak terjelaskan dan gangguan bursa terdokumentasi memutus keadaan indikator. Kebijakan \emph{holiday-reset} juga memutus keadaan indikator pada penutupan yang selaras dengan kalender bursa.


\begin{table}[htbp]
\centering
\small
\setlength{\tabcolsep}{4pt}
\caption{Hasil audit celah harga per sumber}
\label{tab:hasil-audit-celah}
\begin{tabular}{p{7.2cm}rr}
\toprule
Komponen & MetaQuotes & Ultima \\
\midrule
Bar setelah batas awal & 38.799 & 23.135 \\
Seluruh interval celah & 1.705 & 1.010 \\
Penutupan terjadwal & 1.627 & 965 \\
Selaras dengan kalender bursa & 69 & 44 \\
Celah tidak terjelaskan & 9 & 0 \\
Gangguan bursa terdokumentasi & 0 & 1 \\
Reset/segmen kebijakan primer & 9/10 & 1/2 \\
Reset/segmen \emph{holiday-reset} & 78/79 & 45/46 \\
\bottomrule
\end{tabular}
\end{table}


\subsection{Kelayakan Fitur, Target, dan Masker Bersama}

Sampel akhir mensyaratkan fitur yang telah melewati pemanasan, target satu bar yang sah, dan label tiga kelas yang tersedia. Tabel \ref{tab:hasil-kelayakan} menunjukkan bahwa irisan kelayakan fitur dan target pada kebijakan primer berjumlah 36.658 bar untuk MetaQuotes dan 22.031 bar untuk Ultima. Jumlah tersebut mencakup seluruh rentang data setelah penyaringan, termasuk riwayat latih; jumlah observasi uji per periode dilaporkan secara terpisah. Kebijakan \emph{holiday-reset} mengurangi sampel karena pemanasan indikator dimulai kembali pada lebih banyak celah. Bar harga tidak diinterpolasi.


\begin{table}[htbp]
\centering
\small
\setlength{\tabcolsep}{4pt}
\caption{Jumlah observasi yang memenuhi aturan kelayakan pada ambang utama}
\label{tab:hasil-kelayakan}
\begin{tabular}{llrrrr}
\toprule
Sumber & Kebijakan & Input & Fitur & Target & Irisan \\
\midrule
MetaQuotes & Primer & 38.799 & 38.342 & 36.872 & 36.658 \\
MetaQuotes & Holiday-reset & 38.799 & 35.144 & 35.318 & 33.592 \\
Ultima & Primer & 23.135 & 23.037 & 22.078 & 22.031 \\
Ultima & Holiday-reset & 23.135 & 21.052 & 21.112 & 20.128 \\
\bottomrule
\end{tabular}
\end{table}


Distribusi label pada seluruh sampel irisan ditunjukkan pada Tabel \ref{tab:hasil-distribusi-ambang}. Untuk sumber dan kebijakan yang sama, perubahan pengali ambang mengubah label tanpa mengubah jumlah bar yang memenuhi masker. Pada data utama kebijakan primer, jumlah HOLD berubah dari 9.598 menjadi 17.637 dan 27.442 ketika pengali meningkat dari 0,25 menjadi 0,50 dan 1,00. Distribusi ini dilaporkan sebagai sifat target yang dibentuk oleh protokol, bukan sebagai dasar memilih ambang berdasarkan skor model.


\begingroup
\small
\setlength{\tabcolsep}{4pt}
\begin{longtable}{llrrrrr}
\caption{Distribusi label pada seluruh sampel irisan untuk setiap pengali ambang}\label{tab:hasil-distribusi-ambang}\\
\toprule
Sumber & Kebijakan & $\alpha$ & BUY & HOLD & SELL & Jumlah \\
\midrule
\endfirsthead
\caption[]{Distribusi label pada seluruh sampel irisan untuk setiap pengali ambang (lanjutan)}\\
\toprule
Sumber & Kebijakan & $\alpha$ & BUY & HOLD & SELL & Jumlah \\
\midrule
\endhead
\midrule
\multicolumn{7}{r}{Bersambung ke halaman berikutnya} \\
\endfoot
\endlastfoot
MetaQuotes & Primer & 0,25 & 13.873 & 9.598 & 13.187 & 36.658 \\
MetaQuotes & Primer & 0,50 & 9.841 & 17.637 & 9.180 & 36.658 \\
MetaQuotes & Primer & 1,00 & 4.697 & 27.442 & 4.519 & 36.658 \\
MetaQuotes & Holiday-reset & 0,25 & 12.672 & 8.880 & 12.040 & 33.592 \\
MetaQuotes & Holiday-reset & 0,50 & 8.965 & 16.257 & 8.370 & 33.592 \\
MetaQuotes & Holiday-reset & 1,00 & 4.260 & 25.215 & 4.117 & 33.592 \\
Ultima & Primer & 0,25 & 8.297 & 5.867 & 7.867 & 22.031 \\
Ultima & Primer & 0,50 & 5.959 & 10.599 & 5.473 & 22.031 \\
Ultima & Primer & 1,00 & 2.847 & 16.475 & 2.709 & 22.031 \\
Ultima & Holiday-reset & 0,25 & 7.543 & 5.423 & 7.162 & 20.128 \\
Ultima & Holiday-reset & 0,50 & 5.394 & 9.777 & 4.957 & 20.128 \\
Ultima & Holiday-reset & 1,00 & 2.572 & 15.106 & 2.450 & 20.128 \\
\bottomrule
\end{longtable}
\endgroup


\subsection{Sampel Setiap Periode Uji}

Tabel \ref{tab:hasil-sampel-periode} memperlihatkan jumlah observasi latih setelah \emph{purge}, jumlah observasi uji, dan distribusi kelas uji untuk kebijakan primer serta ambang utama. Seluruh periode memiliki tiga kelas benar pada sampel uji. Pada MetaQuotes, ketujuh periode memenuhi persyaratan validasi internal. Total observasi uji MetaQuotes adalah 19.898 bar, dengan setiap bar dipakai bersama oleh K1, K2, K3, serta kedua model. Jumlah latih meningkat seiring bergesernya periode uji; jumlah latih tidak harus sama dengan penjumlahan sederhana jumlah uji yang ditampilkan karena kelayakan horizon label diperiksa ulang pada setiap batas periode.

Pada Ultima, tersedianya bar uji pada tiga periode awal belum mencukupi persyaratan pelatihan: riwayat sebelum September 2022 tidak tersedia untuk membentuk tiga blok validasi internal dalam 18 bulan terakhir. Oleh karena itu, K1, K2, dan K3 pada 2023-H1, 2023-H2, dan 2024-H1 berstatus \textbf{NOT TESTED} untuk kedua model. Empat periode dari 2024-H2 sampai Januari--Juli 2026 dapat dievaluasi. Periode yang gagal tetap menjadi bagian dari tujuh periode terjadwal dan tidak diberi skor rekaan.


\begingroup
\small
\setlength{\tabcolsep}{4pt}
\begin{longtable}{llrrrrr}
\caption{Sampel kebijakan primer dan ambang utama per periode uji}\label{tab:hasil-sampel-periode}\\
\toprule
Sumber & Periode & Latih & Uji & BUY & HOLD & SELL \\
\midrule
\endfirsthead
\caption[]{Sampel kebijakan primer dan ambang utama per periode uji (lanjutan)}\\
\toprule
Sumber & Periode & Latih & Uji & BUY & HOLD & SELL \\
\midrule
\endhead
\midrule
\multicolumn{7}{r}{Bersambung ke halaman berikutnya} \\
\endfoot
\endlastfoot
MetaQuotes & 2023-H1 & 16.756 & 2.807 & 699 & 1.410 & 698 \\
MetaQuotes & 2023-H2 & 19.563 & 2.829 & 669 & 1.470 & 690 \\
MetaQuotes & 2024-H1 & 22.393 & 2.808 & 740 & 1.385 & 683 \\
MetaQuotes & 2024-H2 & 25.201 & 2.870 & 830 & 1.380 & 660 \\
MetaQuotes & 2025-H1 & 28.071 & 2.735 & 816 & 1.229 & 690 \\
MetaQuotes & 2025-H2 & 30.807 & 2.746 & 838 & 1.219 & 689 \\
MetaQuotes & 2026 Jan--Jul & 33.554 & 3.103 & 809 & 1.501 & 793 \\
Ultima & 2023-H1 & 1.834 & 2.807 & 698 & 1.408 & 701 \\
Ultima & 2023-H2 & 4.641 & 2.829 & 673 & 1.467 & 689 \\
Ultima & 2024-H1 & 7.471 & 2.806 & 734 & 1.390 & 682 \\
Ultima & 2024-H2 & 10.277 & 2.868 & 825 & 1.389 & 654 \\
Ultima & 2025-H1 & 13.145 & 2.784 & 835 & 1.252 & 697 \\
Ultima & 2025-H2 & 15.930 & 2.813 & 851 & 1.246 & 716 \\
Ultima & 2026 Jan--Jul & 18.744 & 3.285 & 852 & 1.569 & 864 \\
\bottomrule
\end{longtable}
\endgroup


\section{Hasil Audit Kalender Makroekonomi}

Snapshot kalender ekonomi berisi 706 peristiwa USD pada Januari 2020--Juli 2026 tanpa duplikasi pengenal nilai. Sebanyak 520 baris mempunyai nilai aktual dan 460 baris mempunyai prakiraan. Rinciannya adalah 310 peristiwa CPI, 78 NFP, 78 tingkat pengangguran, 54 keputusan suku bunga, dan 186 peristiwa FOMC. Empat jenis CPI dipertahankan sebagai keluarga fitur terpisah. Dengan NFP dan tingkat pengangguran, terdapat enam keluarga numerik untuk kejutan aktual dikurangi prakiraan. Keputusan suku bunga dan FOMC hanya menyumbang representasi kejadian pascarilis karena prakiraan tidak tersedia.

Pemeriksaan hash memastikan snapshot kalender sama dengan arsip beku. Penandaan UTC khusus snapshot ini cocok dengan sepuluh jangkar rilis resmi lintas musim DST dan non-DST. Pada tujuh periode utama, keenam keluarga numerik mempunyai cakupan nilai kejutan pada data latih dan uji. Penggabungan kalender menggunakan rilis yang waktunya lebih awal secara ketat daripada waktu keputusan bar; rilis tepat pada waktu keputusan baru tersedia bagi keputusan berikutnya. Nilai makro yang tidak tersedia diproses dengan parameter imputasi yang dipelajari hanya dari data latih.

Hasil pemeriksaan tersebut memungkinkan K3 dilaksanakan sesuai asumsi operasional Bab 3. Riwayat independen mengenai nilai aktual rilis pertama dan prakiraan pra-rilis tidak tersedia. Karena itu, hasil K3 dalam bab ini berlaku untuk snapshot retrospektif yang dibekukan dan tidak diklaim sebagai arsip \emph{point-in-time} penuh. Validasi timestamp tidak menghapus keterbatasan versi nilai historis tersebut.

\section{Hasil Pemodelan pada Data Utama}

\subsection{Kinerja K1, K2, dan K3}

Tabel \ref{tab:hasil-logistik} dan Tabel \ref{tab:hasil-xgboost} menyajikan kinerja setiap konfigurasi pada data utama. K1 memakai sembilan fitur harga, K2 menambahkan RSI dan tiga fitur MACD, sedangkan K3 menambahkan representasi kalender. Setiap konfigurasi memilih hiperparameter sendiri melalui validasi internal. Seluruh model pada ketujuh periode utama berhasil dievaluasi; rekaman pencarian pada run utama tidak mencatat kegagalan kandidat dan tidak membuang kolom yang seluruhnya hilang pada data latih.

Pada Regresi Logistik, rerata macro-F1 K1, K2, dan K3 berturut-turut adalah 0,392300, 0,392374, dan 0,390922. Pada XGBoost, reratanya adalah 0,390178, 0,393466, dan 0,393791. Rerata balanced accuracy masing-masing konfigurasi juga ditampilkan agar perubahan macro-F1 dapat dibaca bersama metrik yang merata-ratakan recall tiga kelas. Semua balanced accuracy pada data utama terdefinisi karena setiap periode uji memiliki ketiga kelas benar.


\begin{table}[htbp]
\centering
\small
\setlength{\tabcolsep}{4pt}
\caption{Kinerja Regresi Logistik Multinomial pada data utama}
\label{tab:hasil-logistik}
\begin{tabular}{lrrrrrr}
\toprule
Periode & \multicolumn{3}{c}{Macro-F1} & \multicolumn{3}{c}{Balanced accuracy} \\
 & K1 & K2 & K3 & K1 & K2 & K3 \\
\midrule
2023-H1 & 0,423792 & 0,423555 & 0,413261 & 0,423675 & 0,423698 & 0,414284 \\
2023-H2 & 0,391227 & 0,392407 & 0,389339 & 0,399447 & 0,399807 & 0,399828 \\
2024-H1 & 0,394974 & 0,398701 & 0,395091 & 0,396377 & 0,400367 & 0,398366 \\
2024-H2 & 0,389128 & 0,388312 & 0,387069 & 0,391359 & 0,391458 & 0,390822 \\
2025-H1 & 0,391241 & 0,387247 & 0,389463 & 0,393950 & 0,390010 & 0,392824 \\
2025-H2 & 0,379496 & 0,378420 & 0,382411 & 0,383754 & 0,382757 & 0,386108 \\
2026 Jan--Jul & 0,376240 & 0,377973 & 0,379819 & 0,381888 & 0,383740 & 0,385517 \\
\midrule
Rerata & 0,392300 & 0,392374 & 0,390922 & 0,395778 & 0,395977 & 0,395393 \\
\bottomrule
\end{tabular}
\end{table}


\begin{table}[htbp]
\centering
\small
\setlength{\tabcolsep}{4pt}
\caption{Kinerja XGBoost pada data utama}
\label{tab:hasil-xgboost}
\begin{tabular}{lrrrrrr}
\toprule
Periode & \multicolumn{3}{c}{Macro-F1} & \multicolumn{3}{c}{Balanced accuracy} \\
 & K1 & K2 & K3 & K1 & K2 & K3 \\
\midrule
2023-H1 & 0,424583 & 0,413318 & 0,427637 & 0,425221 & 0,413832 & 0,428617 \\
2023-H2 & 0,393049 & 0,384696 & 0,388074 & 0,409703 & 0,402181 & 0,410556 \\
2024-H1 & 0,392897 & 0,400674 & 0,394101 & 0,398129 & 0,403674 & 0,397388 \\
2024-H2 & 0,387431 & 0,400350 & 0,391979 & 0,389712 & 0,404936 & 0,396001 \\
2025-H1 & 0,389511 & 0,386929 & 0,391868 & 0,390105 & 0,387549 & 0,392325 \\
2025-H2 & 0,376530 & 0,382692 & 0,379501 & 0,377072 & 0,382885 & 0,379592 \\
2026 Jan--Jul & 0,367244 & 0,385605 & 0,383374 & 0,369239 & 0,387055 & 0,385075 \\
\midrule
Rerata & 0,390178 & 0,393466 & 0,393791 & 0,394169 & 0,397444 & 0,398508 \\
\bottomrule
\end{tabular}
\end{table}


\subsection{Diagnostik Prediktor Kelas Mayoritas}

Kelas mayoritas data latih adalah HOLD pada setiap periode utama. Tabel \ref{tab:hasil-mayoritas} menunjukkan kinerja prediktor konstan yang selalu memilih kelas tersebut pada masker uji yang sama. Balanced accuracy diagnostik ini adalah 0,333333 pada seluruh periode, sedangkan macro-F1 berubah menurut proporsi HOLD pada data uji. Macro-F1 dan balanced accuracy K1, K2, maupun K3 pada kedua model lebih tinggi daripada diagnostik ini di masing-masing periode. Perbandingan tersebut menjadi pemeriksaan kewajaran keluaran klasifikasi; pertanyaan kontribusi fitur tetap dijawab melalui pasangan K2--K1 dan K3--K2.


\begin{table}[htbp]
\centering
\small
\setlength{\tabcolsep}{4pt}
\caption{Diagnostik prediktor kelas mayoritas pada data utama}
\label{tab:hasil-mayoritas}
\begin{tabular}{lcrr}
\toprule
Periode & Kelas & Macro-F1 & Balanced accuracy \\
\midrule
2023-H1 & HOLD & 0,222907 & 0,333333 \\
2023-H2 & HOLD & 0,227960 & 0,333333 \\
2024-H1 & HOLD & 0,220208 & 0,333333 \\
2024-H2 & HOLD & 0,216471 & 0,333333 \\
2025-H1 & HOLD & 0,206694 & 0,333333 \\
2025-H2 & HOLD & 0,204960 & 0,333333 \\
2026 Jan--Jul & HOLD & 0,217347 & 0,333333 \\
\bottomrule
\end{tabular}
\end{table}


\subsection{Matriks Kebingungan Gabungan}

Tabel \ref{tab:hasil-matriks-gabungan} menjumlahkan matriks kebingungan tujuh periode uji utama untuk memperlihatkan jumlah prediksi menurut kelas benar dan kelas prediksi. Setiap konfigurasi/model mencakup 19.898 observasi uji yang sama. Baris mengikuti kelas benar dan kolom mengikuti kelas prediksi, dengan urutan BUY, HOLD, SELL. Matriks gabungan digunakan sebagai diagnostik pola kesalahan; rerata macro-F1 dan balanced accuracy pada tabel kinerja tetap dihitung per periode, bukan dihitung ulang dari matriks gabungan.


\begingroup
\small
\setlength{\tabcolsep}{4pt}
\begin{longtable}{lllrrr}
\caption{Matriks kebingungan gabungan tujuh periode data utama}\label{tab:hasil-matriks-gabungan}\\
\toprule
Model & Konfigurasi & Kelas benar & Pred. BUY & Pred. HOLD & Pred. SELL \\
\midrule
\endfirsthead
\caption[]{Matriks kebingungan gabungan tujuh periode data utama (lanjutan)}\\
\toprule
Model & Konfigurasi & Kelas benar & Pred. BUY & Pred. HOLD & Pred. SELL \\
\midrule
\endhead
\midrule
\multicolumn{6}{r}{Bersambung ke halaman berikutnya} \\
\endfoot
\endlastfoot
Logistik & K1 & BUY & 1.448 & 2.389 & 1.564 \\
Logistik & K1 & HOLD & 1.671 & 5.769 & 2.154 \\
Logistik & K1 & SELL & 1.267 & 2.094 & 1.542 \\
Logistik & K2 & BUY & 1.465 & 2.368 & 1.568 \\
Logistik & K2 & HOLD & 1.697 & 5.769 & 2.128 \\
Logistik & K2 & SELL & 1.285 & 2.092 & 1.526 \\
Logistik & K3 & BUY & 1.454 & 2.234 & 1.713 \\
Logistik & K3 & HOLD & 1.772 & 5.521 & 2.301 \\
Logistik & K3 & SELL & 1.308 & 1.942 & 1.653 \\
XGBoost & K1 & BUY & 1.587 & 2.019 & 1.795 \\
XGBoost & K1 & HOLD & 2.065 & 5.118 & 2.411 \\
XGBoost & K1 & SELL & 1.354 & 1.816 & 1.733 \\
XGBoost & K2 & BUY & 1.676 & 1.982 & 1.743 \\
XGBoost & K2 & HOLD & 2.129 & 5.035 & 2.430 \\
XGBoost & K2 & SELL & 1.395 & 1.761 & 1.747 \\
XGBoost & K3 & BUY & 1.663 & 1.943 & 1.795 \\
XGBoost & K3 & HOLD & 2.079 & 5.017 & 2.498 \\
XGBoost & K3 & SELL & 1.396 & 1.725 & 1.782 \\
\bottomrule
\end{longtable}
\endgroup


\subsection{Hiperparameter Terpilih}

Pilihan parameter pada Tabel \ref{tab:hasil-parameter-logistik} dan Tabel \ref{tab:hasil-parameter-xgboost} berasal dari macro-F1 validasi internal, sebelum prediksi periode uji dihitung. Parameter dapat berbeda antarkonfigurasi pada periode yang sama. Oleh sebab itu, perbandingan utama mencakup penambahan kelompok fitur beserta penyetelan ulang model. Pada tabel XGBoost, setiap sel ditulis sebagai $(d,n,\eta)$, yaitu kedalaman maksimum, jumlah iterasi \emph{boosting}, dan laju pembelajaran.


\begin{table}[htbp]
\centering
\small
\setlength{\tabcolsep}{4pt}
\caption{Parameter C terpilih Regresi Logistik pada data utama}
\label{tab:hasil-parameter-logistik}
\begin{tabular}{lrrr}
\toprule
Periode & K1 & K2 & K3 \\
\midrule
2023-H1 & 0,01 & 0,01 & 10 \\
2023-H2 & 1 & 0,01 & 0,01 \\
2024-H1 & 0,01 & 0,01 & 0,1 \\
2024-H2 & 0,1 & 0,01 & 0,1 \\
2025-H1 & 0,1 & 0,01 & 10 \\
2025-H2 & 0,1 & 0,01 & 10 \\
2026 Jan--Jul & 1 & 1 & 0,01 \\
\bottomrule
\end{tabular}
\end{table}


\begin{table}[htbp]
\centering
\small
\setlength{\tabcolsep}{4pt}
\caption{Parameter $(d,n,\eta)$ terpilih XGBoost pada data utama}
\label{tab:hasil-parameter-xgboost}
\begin{tabular}{lccc}
\toprule
Periode & K1 & K2 & K3 \\
\midrule
2023-H1 & $( 3; 100; 0{,}05)$ & $( 3; 100; 0{,}05)$ & $( 3; 100; 0{,}05)$ \\
2023-H2 & $( 3; 100; 0{,}05)$ & $( 3; 100; 0{,}05)$ & $( 3; 100; 0{,}05)$ \\
2024-H1 & $( 3; 100; 0{,}05)$ & $( 5; 100; 0{,}1)$ & $( 3; 300; 0{,}05)$ \\
2024-H2 & $( 3; 100; 0{,}1)$ & $( 5; 300; 0{,}05)$ & $( 5; 300; 0{,}05)$ \\
2025-H1 & $( 3; 300; 0{,}05)$ & $( 5; 300; 0{,}05)$ & $( 5; 300; 0{,}05)$ \\
2025-H2 & $( 3; 100; 0{,}05)$ & $( 5; 300; 0{,}05)$ & $( 3; 300; 0{,}05)$ \\
2026 Jan--Jul & $( 3; 100; 0{,}05)$ & $( 5; 300; 0{,}05)$ & $( 3; 300; 0{,}05)$ \\
\bottomrule
\end{tabular}
\end{table}


\section{Perubahan Kinerja Berpasangan}

\subsection{Penambahan Fitur Teknikal dan Makroekonomi}

Tabel \ref{tab:hasil-paired-f1} menyajikan rata-rata, median, simpangan baku sampel, serta jumlah periode dengan perubahan macro-F1 positif. Semua ringkasan utama memiliki tujuh pasangan dari tujuh periode terjadwal. Regresi Logistik memperoleh rerata $\Delta(K2-K1)$ sebesar +0,000074 dengan tiga periode positif; XGBoost memperoleh +0,003289 dengan empat periode positif. Untuk $K3-K2$, rerata Regresi Logistik adalah $-0{,}001452$ dan XGBoost +0,000324; keduanya memiliki tiga periode positif. Nilai median K3--K2 pada kedua model bernilai negatif.


\begin{table}[htbp]
\centering
\small
\setlength{\tabcolsep}{4pt}
\caption{Ringkasan perubahan macro-F1 pada tujuh pasangan periode utama}
\label{tab:hasil-paired-f1}
\begin{tabular}{llrrrc}
\toprule
Model & Pasangan & Rerata & Median & SD sampel & Positif \\
\midrule
Logistik & $K2-K1$ & +0,000074 & -0,000237 & 0,002452 & 3/7 \\
Logistik & $K3-K2$ & -0,001452 & -0,001243 & 0,004821 & 3/7 \\
XGBoost & $K2-K1$ & +0,003289 & +0,006163 & 0,011031 & 4/7 \\
XGBoost & $K3-K2$ & +0,000324 & -0,002231 & 0,007837 & 3/7 \\
\bottomrule
\end{tabular}
\end{table}


Tabel \ref{tab:hasil-paired-ba} memberikan ringkasan metrik pendamping. Pada K2--K1, Regresi Logistik memiliki lima periode dengan balanced accuracy meningkat, berbeda dari tiga periode dengan macro-F1 meningkat. Pada K3--K2, balanced accuracy Regresi Logistik meningkat pada empat periode, dengan rerata tetap negatif. XGBoost memiliki empat periode positif pada K2--K1 dan tiga pada K3--K2 untuk kedua metrik. Perbedaan ini menunjukkan bahwa arah perubahan yang dicatat harus selalu menyebut metriknya.


\begin{table}[htbp]
\centering
\small
\setlength{\tabcolsep}{4pt}
\caption{Ringkasan perubahan balanced accuracy pada tujuh pasangan periode utama}
\label{tab:hasil-paired-ba}
\begin{tabular}{llrrrc}
\toprule
Model & Pasangan & Rerata & Median & SD sampel & Positif \\
\midrule
Logistik & $K2-K1$ & +0,000199 & +0,000099 & 0,002440 & 5/7 \\
Logistik & $K3-K2$ & -0,000584 & +0,000021 & 0,004339 & 4/7 \\
XGBoost & $K2-K1$ & +0,003276 & +0,005545 & 0,011041 & 4/7 \\
XGBoost & $K3-K2$ & +0,001063 & -0,001980 & 0,008542 & 3/7 \\
\bottomrule
\end{tabular}
\end{table}


Gambar \ref{fig:hasil-delta-periode} memperlihatkan perubahan macro-F1 pada setiap periode utama untuk kedua kelompok fitur. Garis nol memisahkan perubahan positif dan negatif; titik merupakan pasangan skor pada periode yang sama, tanpa penghalusan atau pembobotan jumlah observasi.

\begin{figure}[htbp]
\centering
% Derived from the frozen primary summary; no training or score selection.
\begin{tikzpicture}
\begin{axis}[
width=13.8cm,height=5.6cm,
font=\small, title style={yshift=1.8em}, title={(a) Penambahan teknikal: K2--K1},
xlabel={Periode uji},ylabel={$\Delta$ Macro-F1},
ymin=-0.020,ymax=0.022,ytick={-0.02,-0.01,0,0.01,0.02},
xtick={0,1,2,3,4,5,6},
xticklabels={2023-H1,2023-H2,2024-H1,2024-H2,2025-H1,2025-H2,2026 Jan--Jul},
x tick label style={rotate=25,anchor=east,font=\footnotesize},
ymajorgrids=true,grid style={gray!20},
legend style={at={(0.5,1.04)},anchor=south,legend columns=2,draw=none,font=\small},
scaled y ticks=false,y tick label style={/pgf/number format/fixed,/pgf/number format/precision=2},
]
\addplot[gray,dashed,forget plot] coordinates {(-0.3,0) (6.3,0)};
\addplot[color=blue!75!black,mark=*,thick] coordinates {(0,-0.000237308407) (1,0.001179526437) (2,0.003726554034) (3,-0.000815538988) (4,-0.003993683096) (5,-0.001076196564) (6,0.001733377549)};
\addlegendentry{Logistik}
\addplot[color=orange!80!black,mark=square*,thick] coordinates {(0,-0.011264536589) (1,-0.008352686929) (2,0.007776872361) (3,0.012919193679) (4,-0.002581653699) (5,0.006162663646) (6,0.018360635190)};
\addlegendentry{XGBoost}
\end{axis}
\end{tikzpicture}

\par\vspace{0.25cm}

\begin{tikzpicture}
\begin{axis}[
width=13.8cm,height=5.6cm,
font=\small, title style={yshift=1.8em}, title={(b) Penambahan makro: K3--K2},
xlabel={Periode uji},ylabel={$\Delta$ Macro-F1},
ymin=-0.020,ymax=0.022,ytick={-0.02,-0.01,0,0.01,0.02},
xtick={0,1,2,3,4,5,6},
xticklabels={2023-H1,2023-H2,2024-H1,2024-H2,2025-H1,2025-H2,2026 Jan--Jul},
x tick label style={rotate=25,anchor=east,font=\footnotesize},
ymajorgrids=true,grid style={gray!20},
legend style={at={(0.5,1.04)},anchor=south,legend columns=2,draw=none,font=\small},
scaled y ticks=false,y tick label style={/pgf/number format/fixed,/pgf/number format/precision=2},
]
\addplot[gray,dashed,forget plot] coordinates {(-0.3,0) (6.3,0)};
\addplot[color=blue!75!black,mark=*,thick] coordinates {(0,-0.010293657331) (1,-0.003067441448) (2,-0.003609514790) (3,-0.001243463415) (4,0.002215914534) (5,0.003990983680) (6,0.001845835723)};
\addlegendentry{Logistik}
\addplot[color=orange!80!black,mark=square*,thick] coordinates {(0,0.014318816532) (1,0.003377987405) (2,-0.006573671627) (3,-0.008371057397) (4,0.004938371290) (5,-0.003191020754) (6,-0.002230986516)};
\addlegendentry{XGBoost}
\end{axis}
\end{tikzpicture}

\caption{Perubahan macro-F1 per periode pada konfigurasi utama}
\label{fig:hasil-delta-periode}
\end{figure}

\subsection{Sensitivitas Hiperparameter Tetap}

Dalam analisis pada Tabel \ref{tab:hasil-parameter-tetap}, parameter K1 terpilih dipakai kembali pada K2 dan K3 pada periode serta model yang sama. Praproses tetap dipelajari dari data latih untuk konfigurasi terkait. Pada macro-F1, rerata K2--K1 menjadi +0,000285 untuk Regresi Logistik dan $-0{,}001398$ untuk XGBoost. Rerata K3--K2 menjadi $-0{,}001669$ dan +0,003859. Dengan demikian, arah rerata K2--K1 XGBoost pada analisis ini berbeda dari arah rerata perbandingan utama. Kedua analisis dilaporkan terpisah sesuai desain Bab 3.


\begin{table}[htbp]
\centering
\small
\setlength{\tabcolsep}{4pt}
\caption{Perubahan kinerja dengan hiperparameter K1 tetap pada tujuh pasangan utama}
\label{tab:hasil-parameter-tetap}
\begin{tabular}{lllrrrc}
\toprule
Model & Pasangan & Metrik & Rerata & Median & SD sampel & Positif \\
\midrule
Logistik & $K2-K1$ & Macro-F1 & +0,000285 & +0,000546 & 0,002525 & 4/7 \\
Logistik & $K2-K1$ & BA & +0,000430 & +0,000471 & 0,002516 & 6/7 \\
Logistik & $K3-K2$ & Macro-F1 & -0,001669 & -0,001487 & 0,003744 & 3/7 \\
Logistik & $K3-K2$ & BA & -0,000678 & +0,000353 & 0,003387 & 4/7 \\
XGBoost & $K2-K1$ & Macro-F1 & -0,001398 & +0,001720 & 0,006320 & 4/7 \\
XGBoost & $K2-K1$ & BA & -0,001361 & +0,000547 & 0,006196 & 4/7 \\
XGBoost & $K3-K2$ & Macro-F1 & +0,003859 & +0,003469 & 0,007502 & 5/7 \\
XGBoost & $K3-K2$ & BA & +0,004743 & +0,005468 & 0,007674 & 5/7 \\
\bottomrule
\end{tabular}
\end{table}


\section{Hasil Analisis Sensitivitas}

Tabel \ref{tab:hasil-sensitivitas-teknikal} dan Tabel \ref{tab:hasil-sensitivitas-makro} merangkum perubahan macro-F1 untuk seluruh kombinasi sumber, kebijakan celah, dan ambang yang ditetapkan sebelum skor dibaca. Setiap konfigurasi pada kedua tabel memperoleh penyetelan hiperparameter sendiri sebagaimana analisis utama. Kolom ``Pos.'' menyatakan jumlah pasangan dengan selisih positif dibagi jumlah pasangan sah, sedangkan kolom ``Pas./7'' mempertahankan penyebut tujuh periode terjadwal. Singkatan MQ menunjuk MetaQuotes dan U menunjuk Ultima. Kebijakan HR menunjuk \emph{holiday-reset}; P menunjuk primer. Pengali $\alpha$ membentuk ambang $\alpha\sigma_{t,24}$.

Pada sumber MetaQuotes, rerata perubahan K2--K1 Regresi Logistik positif pada kedua kebijakan untuk ambang 0,25 dan 0,50, tetapi pada ambang 1,00 tandanya berbeda antarkebijakan. Untuk XGBoost, tandanya juga berubah menurut ambang dan kebijakan. Pada K3--K2, Regresi Logistik memperoleh rerata negatif pada kelima kombinasi MetaQuotes dan positif pada HR dengan ambang 1,00. XGBoost memperoleh rerata positif pada empat kombinasi MetaQuotes dan negatif pada dua kombinasi. Ringkasan tersebut melaporkan arah yang teramati tanpa mengganti ambang atau kebijakan utama.


\begingroup
\small
\setlength{\tabcolsep}{4pt}
\begin{longtable}{lccrcrcc}
\caption{Sensitivitas rerata perubahan macro-F1 K2--K1}\label{tab:hasil-sensitivitas-teknikal}\\
\toprule
Sumber & Keb. & $\alpha$ & $\Delta$ Logistik & Pos. & $\Delta$ XGBoost & Pos. & Pas./7 \\
\midrule
\endfirsthead
\caption[]{Sensitivitas rerata perubahan macro-F1 K2--K1 (lanjutan)}\\
\toprule
Sumber & Keb. & $\alpha$ & $\Delta$ Logistik & Pos. & $\Delta$ XGBoost & Pos. & Pas./7 \\
\midrule
\endhead
\midrule
\multicolumn{8}{r}{Bersambung ke halaman berikutnya} \\
\endfoot
\endlastfoot
MQ & P & 0,25 & +0,003388 & 6/7 & +0,000844 & 2/7 & 7/7 \\
MQ & P & 0,50 & +0,000074 & 3/7 & +0,003289 & 4/7 & 7/7 \\
MQ & P & 1,00 & +0,000597 & 3/7 & -0,001379 & 3/7 & 7/7 \\
MQ & HR & 0,25 & +0,003886 & 5/7 & -0,004090 & 1/7 & 7/7 \\
MQ & HR & 0,50 & +0,000127 & 5/7 & +0,000366 & 3/7 & 7/7 \\
MQ & HR & 1,00 & -0,000996 & 1/7 & +0,004996 & 5/7 & 7/7 \\
U & P & 0,25 & +0,001906 & 2/4 & +0,007313 & 4/4 & 4/7 \\
U & P & 0,50 & +0,000325 & 3/4 & +0,002690 & 3/4 & 4/7 \\
U & P & 1,00 & -0,000988 & 2/4 & -0,002315 & 2/4 & 4/7 \\
U & HR & 0,25 & +0,000438 & 3/4 & +0,006806 & 2/4 & 4/7 \\
U & HR & 0,50 & -0,001649 & 2/4 & +0,000146 & 3/4 & 4/7 \\
U & HR & 1,00 & -0,000451 & 2/4 & +0,006385 & 3/4 & 4/7 \\
\bottomrule
\end{longtable}
\endgroup


\begingroup
\small
\setlength{\tabcolsep}{4pt}
\begin{longtable}{lccrcrcc}
\caption{Sensitivitas rerata perubahan macro-F1 K3--K2}\label{tab:hasil-sensitivitas-makro}\\
\toprule
Sumber & Keb. & $\alpha$ & $\Delta$ Logistik & Pos. & $\Delta$ XGBoost & Pos. & Pas./7 \\
\midrule
\endfirsthead
\caption[]{Sensitivitas rerata perubahan macro-F1 K3--K2 (lanjutan)}\\
\toprule
Sumber & Keb. & $\alpha$ & $\Delta$ Logistik & Pos. & $\Delta$ XGBoost & Pos. & Pas./7 \\
\midrule
\endhead
\midrule
\multicolumn{8}{r}{Bersambung ke halaman berikutnya} \\
\endfoot
\endlastfoot
MQ & P & 0,25 & -0,000506 & 4/7 & -0,000782 & 3/7 & 7/7 \\
MQ & P & 0,50 & -0,001452 & 3/7 & +0,000324 & 3/7 & 7/7 \\
MQ & P & 1,00 & -0,002087 & 2/7 & +0,006069 & 5/7 & 7/7 \\
MQ & HR & 0,25 & -0,000692 & 3/7 & +0,002305 & 5/7 & 7/7 \\
MQ & HR & 0,50 & -0,000310 & 3/7 & +0,003173 & 4/7 & 7/7 \\
MQ & HR & 1,00 & +0,001001 & 4/7 & -0,001688 & 3/7 & 7/7 \\
U & P & 0,25 & +0,004383 & 3/4 & +0,004089 & 3/4 & 4/7 \\
U & P & 0,50 & +0,004783 & 3/4 & +0,000649 & 2/4 & 4/7 \\
U & P & 1,00 & -0,001029 & 2/4 & +0,007540 & 3/4 & 4/7 \\
U & HR & 0,25 & +0,004950 & 4/4 & +0,002498 & 2/4 & 4/7 \\
U & HR & 0,50 & +0,004385 & 2/4 & +0,004945 & 2/4 & 4/7 \\
U & HR & 1,00 & +0,000507 & 2/4 & -0,000024 & 2/4 & 4/7 \\
\bottomrule
\end{longtable}
\endgroup


Pada sumber Ultima, hasil utama kebijakan primer untuk empat periode sah ditunjukkan pada Tabel \ref{tab:hasil-ultima-rerata}. Rerata macro-F1 K2--K1 positif pada kedua model, sedangkan K3--K2 masing-masing +0,004783 dan +0,000649. Tabel sensitivitas menunjukkan bahwa arah tersebut tidak bertahan pada seluruh pengali ambang dan kebijakan. Hasil Ultima selalu berasal dari empat pasangan sah dari tujuh periode terjadwal, sehingga tidak mempunyai cakupan waktu uji yang sama dengan ringkasan tujuh periode MetaQuotes.


\begin{table}[htbp]
\centering
\small
\setlength{\tabcolsep}{4pt}
\caption{Rerata kinerja Ultima pada kebijakan primer dan ambang utama}
\label{tab:hasil-ultima-rerata}
\begin{tabular}{lcrrc}
\toprule
Model & Konfigurasi & Macro-F1 & Balanced accuracy & Periode sah \\
\midrule
Logistik & K1 & 0,380734 & 0,387797 & 4/7 \\
Logistik & K2 & 0,381058 & 0,388428 & 4/7 \\
Logistik & K3 & 0,385841 & 0,389725 & 4/7 \\
XGBoost & K1 & 0,380839 & 0,382839 & 4/7 \\
XGBoost & K2 & 0,383529 & 0,385571 & 4/7 \\
XGBoost & K3 & 0,384178 & 0,385436 & 4/7 \\
\bottomrule
\end{tabular}
\end{table}


\section{Ketertelusuran dan Batas Hasil yang Dilaporkan}

Konfigurasi, manifes lingkungan, batas periode, dan artefak fitur dibekukan sebelum pembacaan skor utama. Setiap run menyimpan prediksi per bar, label benar, parameter terpilih, kegagalan kandidat, matriks kebingungan, serta checksum. Pemeriksaan ulang pada 2 Oktober 2026 terhadap dua belas run lengkap K1/K2/K3 memastikan checksum prediksi cocok, metrik dapat dihitung ulang dari prediksi, horizon label tidak melewati batas akhir periode uji, dan K1/K2/K3 menggunakan waktu keputusan serta label uji yang sama. Pemeriksaan ini membaca artefak tersimpan tanpa melatih ulang model atau mengganti protokol. Rerun deterministik yang telah tercatat untuk periode 2023-H1 juga menghasilkan prediksi yang identik dengan run lengkap terkait.

Tabel bab ini ditautkan ke ringkasan ber-hash dalam manifes hasil naskah. Hasil utama memiliki tujuh periode sah; sensitivitas Ultima mempertahankan tiga periode awal sebagai \textbf{NOT TESTED}. K3 tetap dibatasi oleh kalender retrospektif tanpa riwayat revisi independen. Ringkasan berpasangan bersifat deskriptif: simpangan baku dan jumlah periode positif menggambarkan variasi hasil yang tersedia, tanpa uji signifikansi yang mengasumsikan periode saling independen.

\printbibliography[heading=bibintoc,title={DAFTAR REFERENSI}]
\end{document}
